\section{MATH-2}\label{math-2}

\begin{quote}
集合 关系
\end{quote}

\subsection{笛卡尔积}\label{ux7b1bux5361ux5c14ux79ef}

定义，多个集合 \(A_1,A_2,...,A_n\) 的\textbf{笛卡尔积}，

\[
A_1 \times A_2 \times ... \times A_n = 
\{ (x_1, x_2, ..., x_n)\mid x_i \in A_i \}
\]

\subsection{欧式空间}\label{ux6b27ux5f0fux7a7aux95f4}

\(\mathbb{R}^n=\{(x_1, x_2, ... x_n)\mid x_i\in\mathbb{R}\}\) 为 \(n\)
维\textbf{欧几里得空间}

设 \(x,y\in \mathbb{R}^n\)

\[
x=(x_1, x_2, ..., x_n)\\
y=(y_1, y_2, ..., y_n)
\]

记

\[
\begin{aligned}
&|x| =\sqrt{\sum_{i=1}^n x_i^2},\\
&x+y = (x_1+y_1, x_2+y_2, ... , x_n+y_n), \\
&x\cdot y=\sum_{i=1}^n x_i y_i
\end{aligned}
\]

则

\[
\begin{align}
|x+y| &\le |x| +|y| \\
|x\cdot y| &\le |x||y|
\end{align}
\]

\subsection{特征函数}\label{ux7279ux5f81ux51fdux6570}

设 \(X\) 为全集，\(A \subset X\)
定义：\(\chi_A:A\rightarrow \mathbb{R}\)

\[
\chi_A(x)=\left\{\begin{aligned} &1, \space  x \in A \\ 
&0, \space x \in A^C \end{aligned} \right.
\]

\(\chi_A\) 为 \(A\) 的特征函数

\textbf{命题 1} 设 \(X\) 是全集，\(A,B\subset X\)，则

\[
\chi_{A \cup B} = \max\{\chi_A, \chi_B\} \\
\chi_{A \cap B} = \min\{\chi_A, \chi_B\}
\]

\subsection{关系}\label{ux5173ux7cfb}

\textbf{定义 1}

\begin{itemize}
\item
  设 \(X, Y\) 为两个集合，\(f \subset X \times Y\)，则 \(f\)
  为（二元）关系

  如果 \((x,y)\in f\), 记 \(xfy\)

  \(I=\{(x,x) \mid x\in X\}\subset X \times X\) 为恒等关系
\item
  设 \(f \subset X \times Y\)

  \(\mathrm{Dom} f=\{ x\in X \mid \exists y\in Y, \text{s.t.}  (x, y) \in f\}\subset X\)，
  \(f\) 的定义域

  \(\mathrm{Rg} f=\{ y\in Y \mid \exists x\in X, \text{s.t.}  (x, y) \in f\}\subset Y\)，
  \(f\) 的值域
\item
  设 \(f\subset X\times Y, g\subset Y \times Z\)

  \(g \circ f=\{(x, z) \mid \exists y\in Y, \text{s.t.}  (x, y)\in f, (y, z)\in g\} \subset X\times Z\)
\item
  设 \(f\subset X\times Y\)

  \(f^{-1}=\{(y,x)\in Y\times X \mid (x, y) \in f\}\subset Y \times X\)
\end{itemize}

\textbf{结论}

\begin{enumerate}
\def\labelenumi{\arabic{enumi}.}
\tightlist
\item
  设 \(f\subset X\times Y\)，则 \((f^{-1})^{-1}=f\)
\item
  设 \(f\subset X\times Y, g\subset Y\times Z, h \subset Z\times W\)，则
  \((h\circ g)\circ f=h \circ (g\circ f)\)
\item
  设 \(f\subset X\times Y, g\subset  Y\times Z\)，则
  \((g\circ f)^{-1}=f^{-1}\circ g^{-1}\)
\end{enumerate}

\textbf{定义 2}

设 \(\sim \; \subset X \times X\)，\(\sim\) 满足

\begin{enumerate}
\def\labelenumi{\arabic{enumi}.}
\tightlist
\item
  \(\forall x \in X, x \sim x\)
\item
  \(x \sim y \Rightarrow y \sim x\)
\item
  \(x \sim y, y \sim z \Rightarrow x \sim z\)
\end{enumerate}

称 \(\sim\) 为 \(X\) 上的\textbf{等价关系}

设 \(\sim\) 是 \(X\) 上的等价关系，对于任意 \(x\in X\)

称 \([x]_\sim =\{y \in X \mid y \sim x\}\) 为 \(\sim\)
的一个\textbf{等价类}

\section{MATH-3}\label{math-3}

\subsection{关系}\label{ux5173ux7cfb-1}

\(\mathrm{Dom}f^{-1}=\mathrm{Rg}f\)

\(\mathrm{Rg}f^{-1}=\mathrm{Dom}f\)

\(\mathrm{Dom}\space g\circ f=\{x\mid \exists y, z, \space \text{s.t.} (x, y)\in f, (y, z)\in g\}\)

\(\mathrm{Rg}\space g \circ f=\{z\mid \exists x,y,\space\text{s.t.}(x,y)\in f, (y, z)\in g\}\)

等价关系

等价类

设 \([x]\) 为 \(x\) 的等价类，记

\[
X/\sim\;=\{[x]\mid x\in X\}
\]

为 \(X\) 的\textbf{商集}

有

\[
X=\bigcup _{x\in X} [x]
\]

设 \(\sim\) 是 \(X\) 上的等价关系，\(x,y\in X\)，则，

\[
[x]=[y]\Longleftrightarrow x \sim y
\]

序关系

定义 设 \(P\) 是一个集合，\(\le \; \subset P \times P\)，\(\le\) 满足，

\begin{enumerate}
\def\labelenumi{\arabic{enumi}.}
\tightlist
\item
  \(\forall x \in P, x \le x\)
\item
  \(\forall x, y\in P, (x\le y \wedge y\le x \to x=y)\)
\item
  \(\forall x, y, z \in P, (x\le y \wedge y \le z \to x \le z)\)
\end{enumerate}

称 \(\le\) 为 \(P\) 上的\textbf{偏序}，称 \((P, \le)\) 为\textbf{偏序集}

例，\(\le\; = \{(x, y)\in \mathbb{R} ^2 \mid y-x\text{是非负数}\}\) 是
\(\mathbb{R}\) 上的偏序关系

\begin{quote}
注：这句话比较怪，有循环论证之嫌
\end{quote}

全序集

定义 设 \((P, \le)\) 为偏序集

如果 \(\forall x, y \in P\)，必有 \(x\le y \vee y \le x\)

称 \(\le\) 为\textbf{全序}，称 \((P, \le)\) 为\textbf{全序集}

\section{MATH-4}\label{math-4}

\subsection{关系}\label{ux5173ux7cfb-2}

\textbf{定义}

设 \((P, \leq)\) 是全序集，如果
\(\forall A \subset P \wedge A \neq \varnothing\)，\(A\) 有最小元，

称 \(\leq\) 是 \(P\) 上的\textbf{良序}，\((P, \leq)\) 是\textbf{良序集}

\subsection{映射}\label{ux6620ux5c04}

\textbf{定义1}

设 \(f\) 是一个关系，如果 \(f\) 满足，

\[
\forall (x_{1}, y_{1}), (x_{2}, y_{2})\in f, (x_{1}=x_{2}\to y_{1}=y_{2})
\]

称 \(f\) 是\textbf{单值}的，否则称 \(f\) 是\textbf{多值}的。

\textbf{定义2}

设 \(f\) 是一个关系，如果 \(f\) 是单值的，称 \(f\)
为\textbf{函数}（\textbf{映射}，\textbf{变换}，\textbf{算子}，\textbf{映照}\ldots\ldots）

设 \(f\) 是一个函数，\(x\in \mathrm{Dom} \; f\)，则
\(\exists ! y, \; \text{s.t.} (x, y) \in f\)

记 \(y=f(x)\)

称 \(f(x)\) 为 \(x\) 的像

\textbf{命题1}

设 \(f,g\) 是两个函数，则 \(f=g\) 当且仅当，

\begin{enumerate}
\def\labelenumi{\arabic{enumi}.}
\tightlist
\item
  \(\mathrm{Dom} \; f=\mathrm{Dom} \; g\)
\item
  \(\forall x \in \mathrm{Dom} \; f \to f(x)=g(x)\)
\end{enumerate}

\textbf{命题2}

设 \(f,g\) 是两个函数，则 \(g \circ f\) 也是函数，且

\begin{enumerate}
\def\labelenumi{\arabic{enumi}.}
\tightlist
\item
  \(\mathrm{Dom} \; g \circ f= f^{-1}(\mathrm{Dom} \; g)\)
\item
  \((g\circ f)(x)=g(f(x))\)
\end{enumerate}

\begin{quote}
设 \(f\) 为函数，\(A,B\) 为两个集合

\(f(A)=\{f(x)\mid x \in A \cap \mathrm{Dom} \; f\}\)

\(f^{-1}(B)=\{x\in \mathrm{Dom}\; f \mid f(x) \in B\}\)
\end{quote}

\textbf{定义3}

设 \(f\) 是一个函数，如果 \(f\) 满足，

\[
\forall (x_{1},y_{1})\in f, (y_{1}=y_{2}\to x_{1}=x_{2})
\]

称 \(f\) 为\textbf{单射}

\textbf{命题3}

设 \(f\) 为单射，则 \(f^{-1}\) 也是单射

\section{MATH-5}\label{math-5}

\subsection{映射}\label{ux6620ux5c04-1}

\textbf{命题3}

设 \(f\) 是单射，则 \(f^{-1}\) 是单射，并且

\begin{enumerate}
\def\labelenumi{\arabic{enumi}.}
\tightlist
\item
  \(\mathrm{Dom} \; f ^{-1}=\mathrm{Rg} \; f, \;\mathrm{Rg} \; f ^{-1}=\mathrm{Dom}\;f\)
\item
  \(f^{-1}(f(x))=x, \; \forall x \in \mathrm{Dom} \;f\)
\item
  \(f(f^{-1}(y))=y, \; \forall y \in \mathrm{Dom} \; f^{-1}\)
\end{enumerate}

\begin{center}\rule{0.5\linewidth}{0.5pt}\end{center}

\textbf{定义1}

设 \(f\) 是一个函数, \(A=\mathrm{Dom} \; f\),
\(B \supset \mathrm{Rg} \; f\)

称 \(f\) 是 \(A\) 到 \(B\) 的映射，记为 \(f:A \to B\)

\textbf{定义2}

设 \(f:A\to B\)，如果 \(f\) 为单射，称 \(f:A\to B\) 为单射

如果 \(\mathrm{Rg} \; f=B\)，称 \(f:A\to B\) 为满射

如果 \(f:A\to B\) 既是单射又是满射，称 \(f:A\to B\) 是双射

\textbf{命题1}

设 \(f:A\to B\) 为双射，则 \(f^{-1}:B\to A\) 为双射，且

\begin{enumerate}
\def\labelenumi{\arabic{enumi}.}
\tightlist
\item
  \(f^{-1} \circ f = I_{A}\)
\item
  \(f \circ f ^{-1} =I_{B}\)
\end{enumerate}

\begin{quote}
\(I_{X}=\{(x,x) \mid x \in X\}\)
\end{quote}

\subsection{计数}\label{ux8ba1ux6570}

\textbf{定义0}

设 \(S\) 是一个集合，若存在整数 \(n\)，以及双射
\(f:\{1, 2, 3, \dots,n\} \to S\)，则

称 \(S\) 为\textbf{有限集}，记 \(|S|=n\)；我们约定 \(\varnothing\)
为有限集，\(|\varnothing|=0\)

\textbf{定义1}

设 \(A, B\) 是两个集合，如果存在单射 \(f:A\to B\)

称 \(A\) 的\textbf{势}小于等于 \(B\) 的\textbf{势}，记为
\(\overline{\overline{A}}\leq \overline{\overline{B}}\)

\textbf{记号}

如果 \(\overline{\overline{A}}\leq \overline{\overline{B}}\) 也记
\(\overline{\overline{B}}\geq \overline{\overline{A}}\)

如果 \(\overline{\overline{A}}\leq \overline{\overline{B}}\) 且
\(\overline{\overline{B}}\leq \overline{\overline{A}}\)，称 \(A,B\)
等势，记为 \(\overline{\overline{A}}=\overline{\overline{B}}\), 否则记为
\(\overline{\overline{A}}\neq \overline{\overline{B}}\)

如果 \(\overline{\overline{A}}\leq \overline{\overline{B}}\) 且
\(\overline{\overline{A}}\neq \overline{\overline{B}}\)，记为
\(\overline{\overline{A}}< \overline{\overline{B}}\) 或
\(\overline{\overline{B}}>\overline{\overline{A}}\)

\textbf{定理1} (Schröder-Bernstein)

设 \(A,B\) 是两个集合，则
\(\overline{\overline{A}}=\overline{\overline{B}}\) 当且仅当存在双射
\(h:A\to B\)

\subsubsection{必要性证明}\label{ux5fc5ux8981ux6027ux8bc1ux660e}

设 \(f:A\to B,\;g:B\to A\) 为单射

定义 \(F:2^A\to 2^A\)

\begin{quote}
ps. \(2^A\) 表示 \(A\) 的子集的集合，\(2^A=\{x\mid x \subset A\}\)
\end{quote}

\(F(S)=A-g(B-f(S)), \; S\in 2^A\)

考察 \((2^A,\subset)\) 偏序集

如果
\(\forall s_{1},s_{2}\in {2}^{A},s_{1}\subset s_{2}\implies F(s_{1})\subset F(s_{2})\)
称 \(F\) 单调增

\textbf{命题1} \(F\) 单调增

\[
\begin{aligned}
s_{1}\subset s_{2} &\implies f(s_{1})\subset f(s_{2}) \\
&\implies B-f(s_{1}) \supset B-f(s_{2}) \\
&\implies g(B-f(s_{1}))\supset g(B-f(s_{2})) \\
&\implies A-g(B-f(s_{1})) \subset A-g(B-f(s_{2}))\\
&\implies F(s_{1})\subset F(s_{2})
\end{aligned}
\]

如果 \(\exists s \in 2^A, \;\text{s.t.} \; F(S)=S\)，称 \(S\) 为 \(F\)
的不动点

\textbf{命题2} \(F\) 有不动点

记 \(\mathscr{A}=\{ S\in 2^A\mid S\subset F(S) \}\)

先证 \(\mathscr{A} \neq \varnothing\)，因为
\(\varnothing \in \mathscr{A}\)

记 \(C=\bigcup_{S\in \mathscr{A}} S\)

证明 \(C\) 是 \(F\) 的不动点

首先证 \(C\subset F(C)\)，因为
\(C=\bigcup_{S\in\mathscr{A}}S \subset \bigcup_{S\in\mathscr{A}}F(S) \subset F(C)\)

再证
\(C \supset F(C)\)，\(C\subset F(C)\implies F(C)\subset F(F(C))\implies F(C)\in \mathscr{A}\implies F(C)\subset C\)

所以 \(F(C)=C\)

所以 \(g(B-f(C))=A-C\)

则构造 \(h\)

\[
h(x)=\left\{
\begin{aligned}
&f(x), &&x \in C\\
&g^{-1}(x), &&x \in A-C
\end{aligned}
\right.
\]

\(h\) 为双射

\section{MATH-6}\label{math-6}

\subsection{计数}\label{ux8ba1ux6570-1}

\textbf{定义1}

设 \(A\neq\varnothing\) 如果 \(\exists n \in \mathbb{N}^*\) 使得

\[
\overline{\overline{A}}=\overline{\overline{\{1,2,3,\dots,n \}}}
\]

称 \(A\) 为\textbf{有限集}，否则为\textbf{无限集}

\begin{quote}
有限集，存在 \(f:\{ 1,2,\dots,n \}\to A\) 为双射
\end{quote}

\textbf{命题1}

设 \(A\) 为有限集，\(B\) 为无限集，则
\(\overline{\overline{A}}<\overline{\overline{\mathbb{N}^*}}\leq\overline{\overline{B}}\)

\begin{quote}
这意味着： - 有限集的势小于无限集的势 - 正整数集是无限集 -
正整数集的势是所有无限集中最小的
\end{quote}

\textbf{证明}

\begin{enumerate}
\def\labelenumi{\arabic{enumi}.}
\item
  \(\overline{\overline{A}}\leq\overline{\overline{\mathbb{N}^*}}\)

  存在双射 \(f:A\to \{ 1,2,\dots, n \}\)，则 \(f:A \to B\) 是单射
\item
  \(\overline{\overline{A}}\neq\overline{\overline{\mathbb{N}^*}}\)

  反证，假设 \(\exists g:A\to\mathbb{N}^*\)
  是双射，\(\exists f:\{ 1,2,\dots,n\}\to A\) 是双射，则
  \(g\circ f:\{ 1,2,\dots n \}\to\mathbb{N}^*\) 是双射，设
  \(N=\max\{\mathrm{Rg} \; g\circ f\}\)，则
  \(N+1\not\in \mathrm{Rg} \; g\circ f\) 但是 \(N+1\in\mathbb{N}^*\)
\item
  定义 \(\{ a_{n} \},n=1,2,\dots,\text{s.t. } a_{n}\in B\) 并且

  \(a_{n}\not\in\{ a_{1},a_{2},\dots,a_{n-1} \}, \forall n=1,2,\dots\)

  因为 \(B\) 为无限集，\(B\neq\varnothing\)，取 \(a_{1}\in B\)

  因为 \(B\) 为无限集，\(B-\{a_{1}\}\neq\varnothing\)，取
  \(a_{2}\in B-\{ a_{1} \}\)，则 \(a_{2}\in B,a_{2}\not\in \{ a_{1} \}\)

  设 \(a_{n}\in B,a_{n}\not\in \{ a_{1},a_{2},\dots,a_{n-1} \}\)

  因为 \(B\)
  为无限集，\(B-\{ a_{1},a_{2},\dots a_{n} \}\neq \varnothing\)，取
  \(a_{n+1} \in B-\{ a_{1},a_{2},\dots,a_{n} \}\)，则
  \(a_{n+1} \in B, a_{n+1} \not\in \{ a_{1},a_{2},\dots,a_{n} \}\)

  设 \(f:\mathbb{N}^*\to B, \; f(n)=a_{n}, \; n=1,2,\dots\) 为单射，所以
  \(\overline{\overline{\mathbb{N}^*}}\leq\overline{\overline{B}}\)
\end{enumerate}

\textbf{定义2}

设 \(B\) 为无限集，如果
\(\overline{\overline{B}}=\overline{\overline{\mathbb{N}^*}}\)，称 \(B\)
为\textbf{无限可数集}

\textbf{定义3}

设 \(A\neq\varnothing\)，如果 \(A\) 是有限集或无限可数集，称 \(A\)
为\textbf{可数集}，否则称 \(A\) 为\textbf{不可数集}

\begin{enumerate}
\def\labelenumi{\arabic{enumi}.}
\tightlist
\item
  \(\mathbb{Q}\) 为可数集，\(\mathbb{R}\) 为不可数集，当然
  \(\mathbb{R}^2\) 不可数
\item
  可数个可数集的并为可数集
\end{enumerate}

\subsection{实数}\label{ux5b9eux6570}

\subsubsection{实数公理}\label{ux5b9eux6570ux516cux7406}

\begin{enumerate}
\def\labelenumi{\arabic{enumi}.}
\tightlist
\item
  加法公理
\item
  乘法公理
\item
  序公理：\((\mathbb{R},\leq)\) 是全序集
\item
  实数的完备性
\end{enumerate}

\textbf{公理4（完备性）}

设 \(A,B\subset\mathbb{R},\;A,B\neq\varnothing\) 如果

\[
x\leq y,\quad\forall x \in A,y\in B
\]

则 \(\exists c\in\mathbb{R},\text{ s.t. }\)

\[
x\leq c\leq y,\quad \forall x \in A,y\in B
\]

公理4说明数轴上的点与实数集是一一对应的

\subsubsection{确界原理}\label{ux786eux754cux539fux7406}

\textbf{定义1}

设 \(A\subset\mathbb{R},A\neq\varnothing\) 如果
\(\exists M\in \mathbb{R},\text{ s.t. }\)

\[
M\geq x,\quad \forall x \in A
\]

称 \(A\) 有\textbf{上界}，称 \(M\) 为 \(A\) 的\textbf{上界}，如果
\(\exists m\in\mathbb{R},\text{ s.t. }\)

\[
m\leq x,\quad \forall x \in A
\]

称 \(A\) 有\textbf{下界}，称 \(m\) 为 \(A\) 的\textbf{下界}

\textbf{定义2}

设 \(A\subset \mathbb{R}, A\neq\varnothing\)，设 \(A\) 有上界，记

\[
E=\{ s \in\mathbb{R}\mid s\geq x,\; \forall x \in A \}
\]

如果 \(E\) 有最小元，记 \(\sup A=\min E\)，称 \(\sup A\) 为 \(A\)
的\textbf{上确界}，设 \(A\) 有下界，记

\[
F=\{ s \in \mathbb{R} \mid s \leq x,\; \forall x \in A \}
\]

如果 \(F\) 有最大元，记 \(\inf A =\max F\)，称 \(\inf A\) 为 \(A\)
的\textbf{下确界}

\begin{quote}
上确界和下确界如果存在则是唯一的
\end{quote}

\textbf{定理1（确界原理）}

设 \(A \subset \mathbb{R}, A\neq\varnothing\)

如果 \(A\) 有上界，则 \(A\) 必有上确界

如果 \(A\) 有下界，则 \(A\) 必有下确界

\textbf{证明}

设 \(A\) 有上界

\(E=\{ s \in \mathbb{R} \mid s\geq x, \; \forall x \in A \}\)

则 \(A\neq \varnothing, E\neq\varnothing\)，并且

\[
x\leq s, \quad \forall x \in A, s \in E
\]

则 \(\exists c \in \mathbb{R}, \text{ s.t. }\)

\[
x\leq c\leq s, \quad \forall x \in A, s \in E
\]

说明 \(c \in E\)，而且 \(c\) 最小，则 \(\sup A=c\)

下确界同理

\textbf{练习}

\begin{quote}
记 \(-A=\{ -x \mid x \in A \}\)
\end{quote}

设 \(A\subset \mathbb{R}, A\neq\varnothing\)

\begin{enumerate}
\def\labelenumi{\arabic{enumi}.}
\tightlist
\item
  如果 \(A\) 有最大元，则 \(-A\) 有最小元，且 \(-\max A=\min(-A)\)
\item
  如果 \(A\) 有上确界，则 \(-A\) 有下确界，且 \(-\sup A=\inf(-A)\)
\end{enumerate}

\section{MATH-7}\label{math-7}

记 \(\mathbb{Z}=\{ 0, 1, -1, 2, -2, \dots \}\)

设 \(m, n \in \mathbb{Z}, \; m\neq n\)，则 \(|m-n|\geq 1\)

如果 \(m>n\) 则 \(m \geq n+1\)

\textbf{命题1}

设 \(A\subset \mathbb{Z}, \; A \neq \varnothing\)

\begin{enumerate}
\def\labelenumi{\arabic{enumi}.}
\tightlist
\item
  如果 \(A\) 有上界，则 \(A\) 有最大元
\item
  如果 \(A\) 有下界，则 \(A\) 有最小元
\end{enumerate}

\textbf{证明}

对于 1

记 \(M=\sup A\) 则 \(\exists n \in A, \text{ s.t. } n > M-\frac{1}{2}\)

下面证 \(n=\max A\)

设 \(m\in A\)，则 \(m\leq M<n+\frac{1}{2}\) 因此 \(m\leq n\)

（否则 \(m>n \implies m\geq n+1 \wedge m<n+\frac{1}{2}\) 矛盾）

\textbf{命题2}

\(\mathbb{Z}\) 既没有上界，也没有下界

证 因为 \(\mathbb{Z}\) 没有最大元和最小元

\textbf{命题3}（Archimedean）

设 \(h>0\) 则

\[
\mathbb{R} = \bigcup_{k\in\mathbb{Z}} \left[kh, (k+1)h\right)
\]

并且上式右端的集合互不相交

\textbf{推论1}

设 \(x \in \mathbb{R}\) 则
\(\exists !  n \in \mathbb{Z}, \text{ s.t. } n\leq x <n+1\)

称 \(n\) 为 \(x\) 的\textbf{整数部分}

记 \(n=[x]\)

\(x=[x]+(x-[x])\) 其中 \(x-[x]\) 是 \(x\) 的\textbf{小数部分}
\(\in [0, 1)\)

\textbf{推论2}

设 \(x \in \mathbb{R}\)

\begin{enumerate}
\def\labelenumi{\arabic{enumi}.}
\tightlist
\item
  如果 \(x>0\) 则存在
  \(n\in \mathbb{N}^*, \text{ s.t. } 0<\frac{1}{n}<x\)
\item
  如果 \(\forall n \in \mathbb{N}^*,\quad x\leq \frac{1}{n}\)，则
  \(x\leq 0\)
\end{enumerate}

\textbf{命题4}

设 \(a, b\in \mathbb{R},\quad a < b\)

则 \(\exists q \in \mathbb{Q}, \text{ s.t. } a<q<b\)

\(\mathbb{Q} = \left\{  \frac{m}{n} \mid m\in \mathbb{Z}, n \in  \mathbb{N}^*  \right\}\)

\textbf{证明}

\(\exists n \in \mathbb{N}^*, \text{ s.t. } n(b-a)>1\)

则 \(nb>na+1\)

令 \(m=[na]+1\)

\(nb>m>na \implies b> \frac{m}{n} > a\)

\textbf{命题5}

设 \(a, b\in\mathbb{R},\quad a<b\)

则
\(\exists r \in \mathbb{R} \setminus \mathbb{Q}, \text{ s.t. } a<r<b\)

\textbf{证明}

存在 \(q\in\mathbb{Q}, \text{ s.t. } \sqrt{ 2 }a<q<\sqrt{ 2 }b\)

故 \(\frac{q}{\sqrt{ 2 }}\in\mathbb{R} \setminus \mathbb{Q}\) 且
\(a< \frac{q}{\sqrt{ 2 }}<b\)

\section{MATH-8}\label{math-8}

\emph{Inequalities Hardy Littlewood Cambridge}

\subsection{不等式}\label{ux4e0dux7b49ux5f0f}

\(n\) 维欧氏空间

\[
\mathbb{R}^n=\{ (x_{1},x_{2},\dots, x_{n}) \mid x_{i}\in \mathbb{R},\; i=1,2,\dots ,n\}
\]

元素可看成点，向量

设 \(x=(x_{1},x_{2},\dots ,x_{n})\quad y=(y_{1},y_{2},\dots,y_{n})\)

\begin{enumerate}
\def\labelenumi{\arabic{enumi}.}
\tightlist
\item
  加法：\(x+y=(x_{1}+y_{1},x_{2}+y_{2},\dots,x_{n}+y_{n})\)
\item
  数乘：\(\alpha x=(\alpha x_{1}, \alpha x_{2}, \dots , \alpha x_{n}), \quad \alpha \in \mathbb{R}\)
\end{enumerate}

称为向量的\textbf{线性运算}

\(-x=(-1)x \quad x-y=x+(-y)\)

记 \(\mathbf{0}=(0,0,\dots, 0) \in \mathbb{R}^n\)

设 \(\alpha ,\beta \in \mathbb{R}, \quad x, y \in \mathbb{R}^n\)

\(\alpha x+ \beta y\) 是 \(x, y\) 的\textbf{线性组合}

设 \(x, y \in \mathbb{R}^n\) 如果
\(\exists \alpha, \beta \in \mathbb{R}, \alpha^2+\beta^{2}\neq 0\) 使得
\(\alpha x+\beta y = 0 \in \mathbb{R}^n\)

称 \(x, y\) \textbf{线性相关}，否则称 \(x, y\) \textbf{线性无关}（独立）

如果 \(x, y\) 线性相关，又称 \(x \parallel y\) 或共线

\(\mathbf{0} \parallel x, \quad \forall x \in \mathbb{R}^n\)

设
\(x, y \neq 0, \; x \parallel y \Longleftrightarrow \exists k \in \mathbb{R}, \text{ s.t. } y=kx\)

\begin{itemize}
\tightlist
\item
  \(y=kx, \; k>0\)：同向
\item
  \(y=kx,\; k<0\)：反向
\end{itemize}

设 \(a=(a_{1},a_{2},\dots,a_{n}) \quad b=(b_{1}, b_{2},\dots,b_{n})\)

\(a_{i},b_{i}>0, \quad i=1,2,\dots,n\)

\(a\parallel b \Longleftrightarrow \frac{a_{1}}{b_{1}}= \frac{a_{2}}{b_{2}} =\dots=\frac{a_{n}}{b_{n}}\)

设 \(x=(x_{1},x_{2}, \dots,x_{n}) \quad y=(y_{1},y_{2}, \dots, y_{n})\)

定义 \(x \cdot y = \langle x, y \rangle := \sum_{i=1}^n x_{i} y_{i}\)

运算法则

\begin{enumerate}
\def\labelenumi{\arabic{enumi}.}
\tightlist
\item
  \(\forall  x \in \mathbb{R}^n, x\cdot x\geq 0, \text{"="} \Leftrightarrow x=\mathbf{0}\)
\item
  \(\forall x, y \in \mathbb{R}^n,\; x \cdot y =y\cdot x\)
\item
  \(\forall x, y, z \in \mathbb{R}^n, \alpha, \beta \in \mathbb{R},\; \langle \alpha x +\beta y , z\rangle=\alpha \langle x, z\rangle + \beta \langle y, z\rangle\)
  线性性
\end{enumerate}

设 \(x=(x_{1},x_{2},\dots,x_{n})\)

\(|x| =\sqrt{ x \cdot x }=\left( \sum_{i=1}^n  x_{i}^2\right)^{\frac{1}{2}}\)

\(x\) 的长度（模，2-范数）

\begin{enumerate}
\def\labelenumi{\arabic{enumi}.}
\setcounter{enumi}{3}
\tightlist
\item
  \(\forall x \in \mathbb{R}^n, \; |x| \geq 0, \text{"="},\; x=\mathbf{0}\)
\item
  \(\forall \alpha \in \mathbb{ R}, x \in \mathbb{R}^n, \; |\alpha x| = |\alpha| |x|\)
\item
  \(\forall x, y\in \mathbb{R}^n, \; |x+y|\leq|x|+|y|\)
\end{enumerate}

\textbf{命题1}（Cauchy）

设 \(x, y\in \mathbb{R}^n\)，则

\(|x\cdot y |\leq|x| |y|\)

``='' 成立当且仅当 \(x\parallel y\)

\textbf{命题2}（Cauchy）

设 \(a_{i}, b_{i}>0, \; i=1,2,\dots,n\)

\[
\left( \sum_{i=1}^n a_{i}b_{i} \right)^2\leq \left( \sum_{i=1}^n a_{i}^2 \right)\left( \sum_{i=1}^n b_{i} ^{2}\right)
\]

``='' 成立当且仅当
\(\frac{a_{1}}{b_{1}}=\frac{a_{2}}{b_{2}}=\dots=\frac{a_{n}}{b_{n}}\)

拉格朗日恒等式

\$\$

\begin{aligned}

&\left( \sum_{i=1}^n a_{i}^{2} \right)\left( \sum_{i=1}^n b_{i}^2 \right)-\left( \sum_{i=1}^n a_{i}b_{i} \right)^{2}\\
=&\frac{1}{2}\sum_{i=1}^n \sum_{j=1}^n\left(a_{i}b_{j}-a_{j}b_{i}\right)^{2}
\end{aligned}

\$\$

\section{MATH-9}\label{math-9}

\textbf{定理3}（三角不等式）

设 \(x,y\in \mathbb{R}^n\) 则

\[
|x+y|\leq|x|+|y|
\]

设 \(x, y \neq 0\) 则 ``='' 成立，当且仅当 \(x,y\) 同向

证：

\(|x+y|^{2}=|x|^{2}+2x\cdot y+|y|^{2} \leq |x|^{2}+2|x| |y|+|y|^{2}=(|x|+|y|)^{2}\)

\subsection{均值不等式}\label{ux5747ux503cux4e0dux7b49ux5f0f}

\textbf{定理1}（代数、几何平均值不等式 AM-GM）

设 \(a_{i}>0, \; i = 1,2, \dots , n\) 则

\[
\frac{a_{1}+a_{2}+\dots+a_{n}}{n}\geq \sqrt[n]{ a_{1}a_{2}\dots a_{n} }
\]

``='' 成立当且仅当 \(a_{1}=a_{2}=\dots =a_{n}\)

\textbf{定理2}

设
\(\lambda_{i}>0, \; i =1,2,\dots,n, \; \sum_{i=1}^n \lambda_{i}=1, \; a_{i}>0, \; i=1,2,\dots,n\)

则

\[
\lambda_{1}a_{1}+\lambda_{2}a_{2}+\dots+\lambda_{n}a_{n}\geq a_{1}^{\lambda_{1}}a_{2}^{\lambda_{2}}\dots a_{n}^{\lambda_{n}}
\]

``='' 成立当且仅当 \(a_{1}=a_{2}=\dots=a_{n}\)

\textbf{定理3}（Young 不等式）

设 \(p, q>0, \; \frac{1}{p}+\frac{1}{q}=1\) 则

\[
\frac{x^p}{p}+\frac{y^{q}}{q}\geq xy,\quad \forall x, y>0
\]

``='' 成立当且仅当 \(x^{p}=y^{q}\)

设 \(x \in \mathbb{R}^n\)

\[
x = (x_{1},x_{2}, \dots , x_{n})
\]

设 \(p>1\) 定义

\[
|x|_{p}=\left( \sum_{i=1}^n |x_{i}|^{p} \right)^{\frac 1p}
\] \(x\) 的 \(p-\)范数

\begin{enumerate}
\def\labelenumi{\arabic{enumi}.}
\tightlist
\item
  \(|x|_{p}\geq 0\), ``='' 当且仅当 \(x=0\)
\item
  \(|kx|_{p}=|k||x|_{p}, \; \forall k \in \mathbb{R}\)
\item
  \(|x\cdot y|\leq|x|_{p}|y|_{q}, \; p,q > 0, \; \frac{1}{p}+\frac{1}{q}=1\)
\item
  \(|x+y|_{p}\leq |x|_{p}+|y|_{p}\)
\end{enumerate}

\subsection{Hölder不等式}\label{huxf6lderux4e0dux7b49ux5f0f}

\textbf{定理1}（Hölder 不等式）

设 \(\alpha,\beta >0, \; \alpha + \beta =1\) 设
\(a_{i}, b_{i}>0, \; i=1, 2, \dots , n\) 则

\[
\sum_{i=1}^{n}a_{i}^{\alpha}b_{i}^{\beta}\leq \left( \sum_{i=1}^{n}a_{i} \right)^{\alpha}\left( \sum_{i=1}^{n}b_{i} \right)^{\beta}
\]

``='' 成立当且仅当
\(\frac{a_{1}}{b_{1}}=\frac{a_{2}}{b_{2}}=\dots=\frac{a_{n}}{b_{n}}\)

\textbf{定理2}

设 \(r,s \in \mathbb{R}, \; \frac{1}{r}+\frac{1}{s}=1\) 设
\(a_{i},b_{i}>0, \; i=1,2,\dots,n\) 则

\[
\begin{aligned}
\sum_{i=1}^{n}a_{i}b_{i}\leq\left( \sum_{i=1}^{n}a_{i}^{r} \right)^{1/r}\left( \sum_{i=1}^{n}b_{i}^{s} \right)^{1/s}, \quad r>1 \\
\sum_{i=1}^{n}a_{i}b_{i}\geq\left( \sum_{i=1}^{n}a_{i}^{r} \right)^{1/r}\left( \sum_{i=1}^{n}b_{i}^{s} \right)^{1/s}, \quad r<1
\end{aligned}
\]

``='' 成立当且仅当

\[
\frac{a_{1}^{r}}{b_{1}^{s}}=\frac{a_{2}^{r}}{b_{2}^{s}}=\dots=\frac{a_{n}^{r}}{b_{n}^{s}}
\]

\section{MATH-10}\label{math-10}

\subsection{Hölder不等式}\label{huxf6lderux4e0dux7b49ux5f0f-1}

见 {[}\hyperref[math-9]{MATH-9}{]} Hölder不等式定理2

\(\frac{1}{r}+\frac{1}{s}=1\) 称为\textbf{共轭指数}

\(\left( \sum_{i=1}^{n}a_{i}^{r} \right)^{1/r}\) 称为\textbf{幂平均}

Cauchy不等式的推广

\subsection{Minkowski不等式}\label{minkowskiux4e0dux7b49ux5f0f}

\textbf{定理1}（Minkowski）

设 \(r\in \mathbb{R}, \; a_{i}, b_{i} > 0, \; i = 1, 2, \dots , n\) 则

\$\$

\begin{aligned}
\left( \sum_{i=1}^{n} (a_{i}+b_{i})^{r} \right)^{1/r}\leq
\left( \sum_{i=1}^{n} a_{i}^{r} \right)^{1/r}+\left( \sum_{i=1}^{n}b_{i}^{r} \right)^{1/r},  \quad r>1
\\
\left( \sum_{i=1}^{n} (a_{i}+b_{i})^{r} \right)^{1/r}\geq
\left( \sum_{i=1}^{n} a_{i}^{r} \right)^{1/r}+\left( \sum_{i=1}^{n}b_{i}^{r} \right)^{1/r},  \quad r<1

\end{aligned}

\$\$

``='' 成立当且仅当
\(\frac{a_{1}}{b_{1}}=\frac{a_{2}}{b_{2}}=\dots=\frac{a_{n}}{b_{n}}\)

\textbf{证明}（使用Hölder不等式）

\begin{enumerate}
\def\labelenumi{\arabic{enumi}.}
\tightlist
\item
  \(r>1\)（\(r<1\) 同理）
\end{enumerate}

\[
\begin{aligned}
\sum_{i=1}^{n}(a_{i}+b_{i})^{r}&=\sum_{i=1}^{n} (a_{i}+b_{i})^{r-1}(a_{i}+b_{i})
\\
&=\sum_{i=1}^{n}(a_{i}+b_{i})^{r-1}a_{i}+\sum_{i=1}^{n} (a_{i}+b_{i})^{r-1} b_{i}
\\
&\leq \left( \sum_{i=1}^{n}\left((a_{i}+b_{i})^{r-1}\right)^{\frac{r}{r-1}} \right)^{\frac{r-1}{r}}\left[\left( \sum_{i=1}^{n}a_{i}^{r} \right)^{\frac 1r}+\left( \sum_{i=1}^{n}b_{i}^{r} \right)^{\frac 1r} \right]
\\
&=\left( \sum_{i=1}^{n}(a_{i}+b_{i})^{r} \right)^{\frac{r-1}{r}}\left[\left( \sum_{i=1}^{n}a_{i}^{r} \right)^{\frac 1r}+\left( \sum_{i=1}^{n}b_{i}^{r} \right)^{\frac 1r} \right]
\\
\implies
\left( \sum_{i=1}^{n}(a_{i}+b_{i})^{r} \right)^{1/r}&\leq \left( \sum_{i=1}^{n}a_{i}^{r} \right)^{1/r}+\left( \sum_{i=1}^{n}b_{i}^{r} \right)^{1/r}
\end{aligned}
\]

其中放缩一步利用Hölder不等式其中 \(r=\frac{r}{r-1}, \; s=r\)

``='' 成立条件

\[
\left\{
\begin{aligned}
\frac{(a_{1}+b_{1})^{r}}{a_{1}^{r}}=\dots=\frac{(a_{n}+b_{n})^{r}}{a_{n}^{r}}
\\
\frac{(a_{1}+b_{1})^{r}}{b_{1}^{r}}=\dots=\frac{(a_{n}+b_{n})^{r}}{b_{n}^{r}}
\end{aligned}
\right.
\]

然后利用和分比定理即可

\textbf{定理2}（加权的Minkowski不等式）

设
\(r\in \mathbb{R}, \; \alpha_{i}>0, \; \sum_{i=1}^{n}\alpha_{i}=1, \; a_{i},b_{i}>0, \; i=1,2,\dots,n\)

\[
\begin{aligned}
\left( \sum_{i=1}^{n}\alpha_{i}(a_{i}+b_{i})^{r} \right)^{1/r}\leq
\left( \sum_{i=1}^{n} \alpha_{i} a_{i}^{r} \right)^{1/r}+\left( \sum_{i=1}^{n}\alpha _{i}b_{i}^{r} \right)^{1/r}, \quad r>1
\\
\left( \sum_{i=1}^{n}\alpha_{i}(a_{i}+b_{i})^{r} \right)^{1/r}\geq
\left( \sum_{i=1}^{n} \alpha_{i} a_{i}^{r} \right)^{1/r}+\left( \sum_{i=1}^{n}\alpha _{i}b_{i}^{r} \right)^{1/r}, \quad r<1
\end{aligned}
\]

``='' 成立当且仅当
\(\frac{a_{1}}{b_{1}}=\frac{a_{2}}{b_{2}}=\dots=\frac{a_{n}}{b_{n}}\)

\textbf{证明}

\[
\begin{aligned}
\left( \sum_{i=1}^{n}\alpha_{i}(a_{i}+b_{i})^{r} \right)^{1/r}
=&\left( \sum_{i=1}^{n}\left(\alpha_{i}^{1/r}(a_{i}+b_{i})\right)^{r} \right)^{1/r}
\\
\leq&
\left( \sum_{i=1}^{n} \left(\alpha_{i}^{1/r} a_{i}\right)^{r} \right)^{1/r}+\left( \sum_{i=1}^{n}\left(\alpha _{i}^{1/r}b_{i}\right)^{r} \right)^{1/r}
\\
=&\left( \sum_{i=1}^{n} \alpha_{i} a_{i}^{r} \right)^{1/r}+\left( \sum_{i=1}^{n}\alpha _{i}b_{i}^{r} \right)^{1/r}
\end{aligned}
\]

设 \(x=(x_{1},x_{2},\dots,x_{n})\in \mathbb{R}^{n}\)

设 \(1\leq p<\infty\) 定义

\[
|x|_{p}=\left( \sum_{i=1}^{n}|x_{i}|^{p} \right)^{1/p}
\]

称作 \(x\) 的 \(p-\)范数

定义

\[
|x|_{\infty}=\max_{i=1\sim n} |x_{i}|
\]

\(x\) 的 \(\infty-\)范数

\textbf{命题1}

设 \(1\leq p\leq \infty\) 则

\begin{enumerate}
\def\labelenumi{\arabic{enumi}.}
\tightlist
\item
  \(\forall x \in \mathbb{R}^{n}, \; |x|_{p}\geq 0\)，``=''成立当且仅当
  \(x=\mathbf{0}\)
\item
  \(\forall x \in \mathbb{R}^{n}, \;|kx|_{p}=|k||x|_{p}, \; \forall k \in \mathbb{R}\)
\item
  \(\forall x,y \in \mathbb{R}^{n}, \; |x+y|_{p}\leq |x|_{p}+|y|_{p}\)
\item
  \(\forall x, y \in \mathbb{R}^{n}, \; 1<p,q<\infty, \; \frac{1}{p}+\frac{1}{q}=1\)
  或 \(p=1,q=\infty\) 则 \(|x\cdot y|\leq|x|_{p}|y|_{q}\)
\end{enumerate}

\textbf{证明}(4)

设 \(x,y\neq {0}\) 令

\(\hat{x}=\frac{x}{|x|_{p}} \implies |\hat{x}|_{p}=1\)

\(\hat{y}=\frac{y}{|y|_{q}}\implies |\hat{y}|_{q}=1\)

只需证 \(|\hat{x}\cdot\hat{y}|\leq 1\)

\[
\begin{aligned}
|\hat{x}\cdot \hat{y} | &= \left|\sum_{i=1}^{n}\hat{x_{i}} \hat{y_{i}}\right|
\\
&\leq \sum_{i=1}^{n}|\hat{x_{i}}\hat{y_{i}}|
\\
&\leq \sum_{i=1}^{n}\left(\frac{1}{p}|\hat{x_{i}}|^{p}+\frac{1}{q}|\hat{y_{i}}|^{q}\right)
\\
&=\frac{1}{p}+\frac{1}{q}=1
\end{aligned}
\]

\section{MATH-11}\label{math-11}

设 \(1\leq p\leq \infty\) 则

\[
|x+y|_{p}\leq |x|_{p}+|y|_{p}, \quad \forall x, y \in \mathbb{R}^n
\] \textbf{证明}*（对于 \(1<p<\infty\)）

\[
\begin{aligned}
|x+y|_{p}^{p}=&\sum|x_{i}+y_{i}|^{p}
\\
=&\sum|x_{i}+y_{i}|^{p-1}|x_{i}+y_{i}|
\\
\leq&\sum|x_{i}+y_{i}|^{p-1}(|x_{i}|+|y_{i}|)
\\
=&\sum|x_{i}+y_{i}|^{p-1}|x_{i}|+\sum |x_{i}+y_{i}|^{p-1}|y_{i}|
\\
\leq&\left(\sum\left(|x_{i}+y_{i}|^{p-1}\right)^{\frac{p}{p-1}} \right)^{\frac{p-1}{p}}\left[\left( \sum|x_{i}|^{p} \right)^{\frac{1}{p}}+\left( \sum|y_{i}|^{p} \right)^{\frac{1}{p}}\right]
\\
=&|x+y|_{p}^{p-1}\left(|x|_{p}+|y|_{p}\right)
\end{aligned}
\]

\subsection{幂平均值不等式}\label{ux5e42ux5e73ux5747ux503cux4e0dux7b49ux5f0f}

\textbf{定义}

设 \(q_{i}>0, \; i=1\sim n, \; \sum_{i=1}^{n}q_{i}=1\)

设 \(r\in \mathbb{R}, \; r\neq 0\)

设 \(a_{i}>0, \; i=1\sim n, \; a=(a_{1},a_{2},\dots, a_{n})\)

定义

\[
m_{r}(a)=\left( \sum_{i=1}^{n}q_{i}a_{i}^{r} \right)^{\frac{1}{r}}
\]

称为 \(a_{1},a_{2},\dots,a_{n}\) 的\textbf{加权的幂平均值}

幂平均值不等式如下

\[
m_{r}(a)\leq m_{s}(a), \quad r\leq s
\]

令
\(\frac{1}{a}=\left( \frac{1}{a_{1}}, \frac{1}{a_{2}}, \dots, \frac{1}{a_{n}} \right)\)

则有

\[
m_{-r}(a)=\frac{1}{m_{r}\left( \frac{1}{a} \right)}
\]

\textbf{命题1} 当
\(r\to + \infty, \; m_{r}(a)\to \max_{1\leq i\leq n} a_{i}\)

\textbf{证}

\(M=\max_{1\leq i\leq n} a_{i}\)

设 \(a_{\alpha}=M\)

设 \(r>0\)

则 \(m_{r}(a)=\left( \sum_{i=1}^{n}q_{i}a_{i}^{r} \right)^{1/r}\leq M\)

\(m_{r}(a)\geq(q_{\alpha}a_{\alpha}^{r})^{1/r}=q_{\alpha}^{1/r}M\)

故

\[
q_{\alpha}^{1/r}M\leq m_{r}(a)\leq M, \quad \forall r>0
\]

在 \(r\to +\infty, \; q_{\alpha}^{1/r}M\to M\)

所以得证

\textbf{命题2} 当
\(r\to - \infty, \; m_{r}(a)\to \min_{1\leq i\leq n} a_{i}\)

用 \(G(a)=a_{1}^{q_{1}}a_{2}^{q_{2}}\dots a_{n}^{n}\) 几何平均值

\textbf{命题3} 当 \(r\to 0, \; m_{r}(a)\to G(a)\)

\textbf{证}

L'Hôspital 法则

\[
m_{r}(a)=\left( \sum_{i=1}^{n}q_{i}a_{i}^{r} \right)^{1/r}
\]

\[
\lim_{ r \to 0 } \ln m_{r}(a)= \lim_{ r \to 0}  \frac{\ln\left( \sum_{i=1}^{n}q_{i}a_{i}^{r} \right)}{r}  = \lim_{ r \to 0 } \frac{\frac{\sum q_{i} a_{i}^{r} \ln a_{i}}{\sum q_{i}a_{i}^{r}}}{1}=\frac{\sum q_{i} \ln a_{i}}{\sum q_{i}} = \ln G(a)
\] 即证

\textbf{定义}

\(m_{0}(a)=\lim_{ r \to 0 } m_{r}(a)=G(a)\)

\(m_{+\infty}=\lim_{ r \to +\infty }m_{r}(a)= \max a\)

\(m_{-\infty}=\lim_{ r \to -\infty }m_{r}(a)=\min a\)

记

\([-\infty, +\infty]=\mathbb{R} \cup\{-\infty, +\infty  \}\)
\textbf{广义实数集}

\textbf{定理1}

设 \(-\infty\leq r<s\leq +\infty\) 则

\[
m_{r}(a)\leq m_{s}(a)
\]

``='' 成立当且仅当 \(a_{1},a_{2},\dots,a_{n}\) 全相等

\textbf{证明}

step1. 证明如果 \(a_{1},a_{2},\dots,a_{n}\) 全相等，则 ``='' 成立

下面设 \(a_{1},a_{2}, \dots,a_{n}\) 不全相等，证 \(m_r(a)<m_s(a)\)

step2. 设 \(0<r<+\infty\)，证 \(m_{r}(a)<m_{+\infty}(a)\)

step3. 设 \(0<r<s<+\infty\)，证 \(m_{r}(a)<m_{s}(a)\)

设 \(r=\alpha s, \; \alpha \in (0, 1)\)

\[
\begin{aligned}
m_{r}(a)=&\left( \sum q_{i} a_{i}^{r} \right)^{1/r} \\
=& \left( \sum q_{i}a_{i}^{\alpha s} \right)^{1/r}\\
=& \left( \sum (q_{i}a_{i}^{s})^{\alpha} q_{i}^{1-\alpha}\right)^{1/r}\\
<&\left( \left( \sum q_{i}a_{i}^{s} \right)^{\alpha} \left( \sum q_{i}^{1-\alpha} \right)\right)^{1/r}\\
=& \left( \sum q_{i}a_{i}^{s} \right)^{1/s}=m_{s}(a)
\end{aligned}
\]

step4. 设 \(r>0\)，这 \(G(a)=m_{0}(a)<m_{r}(a)\)

使用加权算术几何均值不等式

step5. 设 \(r<0\)，证 \(m_{r}(a)<m_{0}(a)\)

step6. 设 \(-\infty<r<s<0\)，证 \(m_{r}(a)<m_{s}(a)\)

step7. 设 \(-\infty<r<0\)，证 \(m_{-\infty}(a)<m_{r}(a)\)

即证

\section{MATH-13}\label{math-13}

\$\$

\begin{aligned}

a_{1}=\frac{3}{2}\quad
\frac{a_{n}}{n}= \frac{3a_{n-1}}{2a_{n-1}+n-1} \\

\frac{n}{a_{n}}=\frac{2}{3}+\frac{1}{3}\cdot \frac{n-1}{a_{n-1}}\\
b_{n}=\frac{n}{a_{n}}\\
b_{1} =\frac{2}{3} \quad b_{n}=\frac{2}{3}+\frac{b_{n-1}}{3} \\

b_{n}-1= \frac{b_{n-1}-1}{3} \\
b_{n}-1=\frac{1}{3^{n-1}}\cdot \left( -\frac{1}{3} \right)\\
b_{n}=1 - \frac{1}{3^n}\\
a_{n}=\frac{n}{1-\frac{1}{3^n}} \\

\prod_{i=1}^{n}a_{i} =\frac{n!}{\prod_{i=1}^{n}\left(1-\frac{1}{3^i} \right)}< 2 n ! \\
\Leftrightarrow 

\prod_{i=1}^{n}\left(1-\frac{1}{3^{i}}\right) >\frac{1}{2}

\end{aligned}

\$\$

\subsection{排序不等式}\label{ux6392ux5e8fux4e0dux7b49ux5f0f}

\textbf{定理1}（Abel求和变换）

设 \(n\in \mathbb{N}^*, \; n\geq 2\) 设
\(a_{i}, b_{i}\in \mathbb{R}, \; i=1\sim n\)

\[
B_{i}=\sum_{k=1}^{i} b_{k}, \quad i=1\sim n
\]

则

\[
\sum_{i=1}^{n}a_{i}b_{i}=a_{n}B_{n}-\sum_{i=1}^{n-1}(a_{i+1}-a_{i})B_{i}
\]

\textbf{定理2}（排序不等式）

设 \(n \in \mathbb{N}^{*}, \; n\geq 2\) 设
\(a_{i}, b_{i}\in \mathbb{R}, i=1\sim n\)

\[
\begin{aligned}
a_{1}\leq a_{2}\leq\dots\leq a_{n}\\
b_{1}\leq b_{2}\leq\dots\leq b_{n}
\end{aligned}
\]

设 \(\sigma : \{1,2,\dots,n\} \to \{1,2, \dots ,n\}\) 为双射，则

\[
\sum_{i=1}^{n}a_{i}b_{n+1-i}\leq \sum_{i=1}^{n}a_{i}b_{\sigma(i)}\leq \sum_{i=1}^{n}a_{i}b_{i}
\]

并且 \(\sum_{i=1}^{n}a_{i}b_{n+1-i}=\sum_{i=1}^{n}a_{i}b_{i}\)

当且仅当 \(a_{1},\dots,a_{n}\) 全相等或 \(b_{1},\dots,b_{n}\) 全相等

\textbf{证明}

乱序和大于等于逆序和

\[
\begin{aligned}
&\sum_{i=1}^{n}a_{i}b_{\sigma(i)}-\sum_{i=1}^{n}a_{i}b_{n+1-i}\\
=&\sum_{i=1}^{n}a_{i}\left(b_{\sigma(i)}-b_{n+1-i}\right)\\
=&\;a_{n}\left( \sum_{i=1}^{n} b_{\sigma(i)}- \sum_{i=1}^{n}  b_{n+1-i} \right)-\sum_{i=1}^{n-1} \left[(a_{i+1}-a_{i})\left(\sum_{k=1}^{i}b_{\sigma(k)}-\sum_{k=1}^{i} b_{n+1-k} \right)\right]\\
\geq &\; 0
\end{aligned}
\]

正序和大于等于乱序和，取负即可

\textbf{加权均值不等式}

设
\(\alpha_{1}, \alpha_{2}, \dots ,\alpha_{n}>0,\; \sum_{i=1}^{n}\alpha_{i}=1\)
设 \(a_{i}>0, \; i=1\sim n\) 则

\[
\sum_{i=1}^{n}\alpha_{i}a_{i}\geq \prod_{i=1}^{n}a_{i}^{\alpha_{i}}
\]

``='' 成立当且仅当 \(a_{i}\) 全相等

\textbf{证明}

\(f(x)=-\ln x \quad f'(x)=\frac{1}{x^{2}}>0\) 严格凸

\(f\left( \sum_{i=1}^{n }\alpha_{i}a_{i} \right)\leq \sum_{i=1}^{n}\alpha_{i}f(a_{i})\)

``=''成立 \(\Leftrightarrow\) \(a_{i}\) 全相等

\[
\ln\left( \sum_{i=1}^{n}\alpha_{i}a_{i} \right)\geq \sum_{i=1}^{n}\alpha_{i}\ln(a_{i})
=\ln\left( \prod_{i=1}^{n}a_{i}^{\alpha_{i}} \right)
\]

即证

\subsection{切比雪夫不等式}\label{ux5207ux6bd4ux96eaux592bux4e0dux7b49ux5f0f}

\textbf{定理1}（Chebyshev单调不等式）

设 \(q_{i}>0,\;  i=1\sim n, \; \sum_{i=1}^{n}q_{i}=1\) 设
\(a_{i},b_{i}\in \mathbb{R}, \; i=1\sim n\) 如果

\[
\begin{aligned}
a_{1}\leq a_{2}\leq\dots\leq a_{n}\\
b_{1}\leq b_{2}\leq \dots \leq b_{n}
\end{aligned}
\]

则

\[
\sum_{i=1}^{n}q_{i}a_{i}b_{i}\geq\left( \sum_{i=1}^{n}q_{i}a_{i} \right)\left( \sum_{i=1}^{n} q_{i}b_{i} \right)
\]

如果

\[
\begin{aligned}
a_{1}\leq a_{2}\leq\dots\leq a_{n}\\
b_{1}\geq b_{2}\geq\dots\geq b_{n}
\end{aligned}
\]

则

\[
\sum_{i=1}^{n}q_{i}a_{i}b_{i}\leq \left( \sum_{i=1}^{n}q_{i}a_{i} \right)\left( \sum_{i=1}^{n}q_{i}b_{i} \right)
\]

``='' 成立当且仅当 \(a_{i}\) 全相等或 \(b_{i}\) 全相等

\textbf{推论}

如果

\[
\begin{aligned}
a_{1}\leq a_{2}\leq\dots\leq a_{n}\\
b_{1}\leq b_{2}\leq\dots\leq b_{n}
\end{aligned}
\]

则

\[
\sum_{i=1}^{n}a_{i}b_{i}\geq \frac{1}{n}\left( \sum_{i=1}^{n}a_{i} \right)\left( \sum_{i=1}^{n}b_{i} \right)\geq \sum_{i=1}^{n}a_{i}b_{n+1-i}
\]

``='' 成立当且仅当 \(a_i\) 全相等或 \(b_{i}\) 全相等

\textbf{证明}（Chebyshev单调不等式）

注意到以下恒等式

\[
\begin{aligned}
&\sum_{i=1}^{n} \sum_{j=1}^{n} q_{i}q_{j}(a_{i}-a_{j})(b_{i}-b_{j})\\
=& \sum_{i=1}^{n}\sum_{j=1}^{n}(q_{i}q_{j}a_{i}b_{i}+q_{i}q_{j}a_{j}b_{j}-q_{i}q_{j}a_{j}a_{i}-q_{i}q_{j}a_{i}b_{j})\\
=& \;  2\left(\sum_{i=1}^{n}\sum_{j=1}^{n}q_{i}a_{i}b_{i}q_{j} - \sum_{i=1}^{n}\sum_{j=1}^{n}q_{i}b_{i}q_{j}a_{j}\right)\\
=& \;  2\left[\left(\sum_{i=1}^{n}q_{i}a_{i}b_{i}\right) - \left(\sum_{i=1}^{n}q_{i}b_{i}\right)\left(\sum_{i=1}^{n}q_{i}a_{i}\right)\right]\\
\end{aligned}
\]

轻松易证上述不等式

\subsection{伯努利不等式}\label{ux4f2fux52aaux5229ux4e0dux7b49ux5f0f}

\textbf{定理1}

设 \(a\in \mathbb{R}\) 如果 \(a>1\) 则

\[
(1+x)^{a} > 1+ax, \quad \forall \; x> -1, \; x\neq 0
\]

如果 \(0<a<1\)，则

\[
(1+x)^{a}<1+ax, \quad \forall \; x >-1, \; x\neq 0
\]

\section{MATH-14}\label{math-14}

\textbf{推论}

设 \(n\in \mathbb{N}^*, \; n\ge {2}\) 则

\[
(1+x)^{n}>1+nx, \quad \forall\;x>-1, \; x\neq 0
\]

\textbf{定理2}

设
\(n\in \mathbb{N}^{*}, \; n\geq 2, \; x_{i}>-1, \; i=1,2,\dots,n\)，并且
\(x_{i}, \; i=1\sim n\) 同正或同负，则

\[
\prod_{i=1}^{n}(1+x_{i})>1+\sum_{i=1}^{n}x_{i}
\]

\textbf{证明}

数学归纳法

\begin{enumerate}
\def\labelenumi{\arabic{enumi}.}
\tightlist
\item
  \(n=2\)
\end{enumerate}

\((1+x_{1})(1+x_{2})=1+x_{1}+x_{2}+x_{1}x_{2}>1+x_{1}+x_{2}\)

\begin{enumerate}
\def\labelenumi{\arabic{enumi}.}
\setcounter{enumi}{1}
\tightlist
\item
  假设 \(n=k\) 成立
\end{enumerate}

若 \(n=k+1\)

\(\prod_{i=1}^{n} (1+x_{i})>\left( 1+\sum_{i=1}^{n-1}x_{i} \right)(1+x_{n})=1+\sum_{i=1}^{n}x_{i}+x_{n}\sum_{i=1}^{n-1}x_{i}>1+\sum_{i=1}^{n}x_{i}\)

\subsection{函数}\label{ux51fdux6570}

\textbf{定义1}

设 \(D\subset \mathbb{R}, \; D \neq \varnothing\)

\[
\forall x \in D,\quad -x \in D
\]

称 \(D\) 关于原点对称，设 \(f: D\to \mathbb{R}\)，如果

\[
f(-x)=f(x), \quad \forall x \in D
\]

称 \(f\) 为\textbf{偶函数}，如果

\[
f(-x)=-f(x),\quad \forall x \in D
\]

称 \(f\) 为\textbf{奇函数}

\textbf{例1}

\begin{enumerate}
\def\labelenumi{\arabic{enumi}.}
\tightlist
\item
  \(f(x)=\ln\left(\sqrt{ x^{2}+1 }-1\right)\)
\item
  \(f(x)=\ln \frac{1-x}{1+x}\)
\item
  \(f(x)=\frac{x}{e^{x}-1}+\frac{x}{2}\)
\end{enumerate}

偶函数：1,3 奇函数：2

p.s. 我是人机+1

\textbf{例2}

设 \(f:\mathbb{R}\to \mathbb{R}\) 无限可微，\(f^{(k)}\) 为 \(f\) 的
\(k\) 阶导

若 \(f\) 为偶函数，讨论 \(f^{(k)}\) 的奇偶性

p.s. 我是人机+2

\textbf{定义2}

设 \(f:\mathbb{R}\to \mathbb{R}\) 如果存在 \(T\neq 0\), 使得

\[
f(x+T)=f(x), \quad \forall x \in \mathbb{R}
\]

称 \(f\) 为 \(T\)-周期函数

\textbf{例1}

证明

\[
D(x)=
\left\{
\begin{aligned}
1, \quad x为有理数\\
0, \quad x为无理数
\end{aligned}
\right.
\]

为周期函数

\textbf{例2}

存在一个非常值的周期函数 \(f\)

使得 \(1,\sqrt{ 2 }\) 为 \(f\) 的周期

\textbf{选择公理}

设 \(\mathscr{A}\)
是一个集族，\(\mathscr{A}\neq \varnothing, \; \forall A \in \mathscr{A}, \; A \neq \varnothing\)

则
\(\exists \varphi: \mathscr{A} \to \bigcup_{A\in \mathscr{A}}A\)，使得

\[
\varphi(A)\in A, \quad \forall A \in \mathscr{A}
\]

\(\varphi\) 称为\textbf{选择函数}

\textbf{证明}（例2）

定义
\(E = \{ (x, y)\in \mathbb{R}^{2} \mid \exists m,n\in \mathbb{Z}, \text{ s.t. } y-x=m+n\sqrt{ 2 } \}\)

\(E\) 是一个 \(\mathbb{R}\) 上的等价关系

\begin{enumerate}
\def\labelenumi{\arabic{enumi}.}
\tightlist
\item
  \(xEx\) 证 \(x-x=0+0\sqrt{ 2 }\)
\item
  \(xEy\to yEx\) 证 \(y-x=m+n\sqrt{ 2 }\to x-y=-m +(-n)\sqrt{ 2 }\)
\item
  \(xEy \wedge yEz \to xEy\) 证
  \(y-x=m_{1}+n_{1}\sqrt{ 2 } \wedge z-y=m_{2}+n_{2}\sqrt{ 2 }\to z-x=(m_{1}+m_{2})+(n_{1}+n_{2})\sqrt{ 2 }\)
\end{enumerate}

\([x]\) 表示 \(x\) 的等价类

设 \(\varphi:\mathbb{R} / E \to \mathbb{R}\) 为选择函数

\[
\varphi([x])\in [x], \quad \forall x \in \mathbb{R}
\]

定义 \(f:\mathbb{R}\to \mathbb{R}\)

\[
f(x)=\varphi([x]), \quad x \in \mathbb{R}
\]

\begin{enumerate}
\def\labelenumi{\arabic{enumi}.}
\tightlist
\item
  \(f(x)=f(y) \Leftrightarrow xEy\)
\item
  \(f(0)\neq f(\sqrt{3})\)
\item
  \(f(x+1)=f(x)\)
\item
  \(f(x+\sqrt{ 2 })=f(x)\)
\end{enumerate}

其中 2 的证明

反证 \(f(0)=f(\sqrt{ 3 })\)

可得 \(0E\sqrt{ 3 }\)

即
\(\sqrt{ 3 }=m+n\sqrt{ 2 } \implies 3=m^{2}+2n^{2}+2\sqrt{ 2 }mn\implies mn=0\)

\begin{enumerate}
\def\labelenumi{\arabic{enumi}.}
\tightlist
\item
  \(m=0\) 则 \(3=2n^{2}\implies \perp\)
\item
  \(n=0\) 则 \(3=m^{2}\implies \perp\)
\end{enumerate}

\textbf{例3}

设 \(f:\mathbb{R}\to \mathbb{R}\) 证明：\(\exists!\) 偶函数
\(g:\mathbb{R}\to \mathbb{R}\)，奇函数 \(h:\mathbb{R}\to \mathbb{R}\)

使得 \(f=g+h\)

\(g(x)=\frac{f(x)+f(-x)}{2}, \; h(x)=\frac{f(x)-f(-x)}{2}\)

\section{MATH-15}\label{math-15}

\[
\begin{aligned}
f(x)=ae^{x-1}-\ln x+\ln a \geq 1 \\ 
\iff e ^{\ln a+x-1} +\ln a+x-1\geq x+\ln x\\
g(x)=e^{x}+x\\
g(\ln a+x-1)\geq g(\ln x)\\
\ln a+x-1\geq \ln x\\
\ln a\geq \ln x-x+1\\
h(x)=\ln x-x+1\\
\ln a\geq h_{max}(x)\\
h'(x)=\frac{1}{x}-1=\frac{1-x}{x}\\
h_{max}(x)=h(1)=0\\
\ln a\geq 0 \iff a\geq 1
\end{aligned}
\]

\textbf{例1}

考察下列函数的单调性

\begin{enumerate}
\def\labelenumi{\arabic{enumi}.}
\item
  \(f(x)=(1+x)^{\frac{1}{x}}, \quad x>0\)
\item
  \(f(x)=(1+x)^{\frac{1}{x}+1}, \quad x>0\)
\item
\end{enumerate}

\[
\begin{aligned}
&g(x) = \ln f(x) = \frac{1}{x} \cdot \ln (1+x)\\
&g'(x) = \frac{\frac{x}{1+x}-\ln(1+x)}{x^{2}} = \frac{f'(x)}{f(x)}\\
&-\ln(1+x)=\ln\left( 1-\frac{x}{1+x} \right)< -\frac{x}{1+x} \\
&\implies g'(x)<0 \implies f'(x)<0
\end{aligned}
\]

\begin{enumerate}
\def\labelenumi{\arabic{enumi}.}
\setcounter{enumi}{1}
\tightlist
\item
\end{enumerate}

\[
\begin{aligned}
&g(x) = \ln f(x) =\left(1+\frac{1}{x} \right) \cdot \ln (1+x) = \frac{(1+x)\ln(1+x) }{x}\\
&g'(x) = \frac{x+x\ln(1+x)-(1+x)\ln(1+x)}{x^{2}}=\frac{x-\ln(1+x)}{x^{2}}=\frac{f'(x)}{f(x)}\\
&\ln(1+x)<x\implies g'(x)>0\implies f'(x)>0
\end{aligned}
\]

设 \(f:A\subset \mathbb{R }\to \mathbb{R}\) 如果
\(\exists M \in \mathbb{R},\text{ s.t. }\)

\[
f(x)\leq M, \quad \forall x \in A
\]

称 \(f\) 有\textbf{上界}，称 \(M\) 为 \(f\)
的上界，类似可定义\textbf{下界}

如果 \(f\) 既有上界又有下界，称 \(f\) \textbf{有界}

\textbf{例1}

考察下列函数的有界性

\begin{enumerate}
\def\labelenumi{\arabic{enumi}.}
\item
  \(f(x)=x^{\frac{1}{x}}, \quad x>0\)
\item
  \(f(x)=(1+x)^{\frac{1}{x}}, \quad x>0\)
\item
  \(f(x)=(1+x)^{\frac{1}{x}+1}, \quad x>0\)
\item
\end{enumerate}

\[
\begin{aligned}
\ln f(x)=\frac{\ln x}{x}\leq \frac{x-1}{x} < 1
\end{aligned}
\]

有界

\begin{enumerate}
\def\labelenumi{\arabic{enumi}.}
\setcounter{enumi}{1}
\tightlist
\item
\end{enumerate}

\[
\ln f(x)= \frac{\ln(1+x)}{x} \leq \frac{x}{x} =1
\]

有界

\begin{enumerate}
\def\labelenumi{\arabic{enumi}.}
\setcounter{enumi}{2}
\tightlist
\item
\end{enumerate}

\[
f(x)=(1+x)^{1/x+1}\geq 1+x
\]

显然无上界

计算这三个函数的导数

\begin{enumerate}
\def\labelenumi{\arabic{enumi}.}
\tightlist
\item
\end{enumerate}

\[
\begin{aligned}
g(x)=\ln f(x)=\frac{\ln x}{x}\\
g'(x)=\frac{1-\ln x}{x^{2}}=\frac{f'(x)}{f(x)}\\
f'(x)=\left( \frac{1-\ln x}{x^{2}} \right)x^{\frac{1}{x}}
\end{aligned}
\]

\begin{enumerate}
\def\labelenumi{\arabic{enumi}.}
\setcounter{enumi}{1}
\tightlist
\item
\end{enumerate}

\[
\begin{aligned}
g(x) = \ln f(x) = \frac{1}{x} \cdot \ln (1+x)\\
g'(x) = \frac{\frac{x}{1+x}-\ln(1+x)}{x^{2}} = \frac{f'(x)}{f(x)}\\
f'(x)=\left( \frac{\frac{x}{1+x}-\ln(1+x)}{x^{2}} \right)(1+x)^{\frac{1}{x}}
\end{aligned}
\]

\begin{enumerate}
\def\labelenumi{\arabic{enumi}.}
\setcounter{enumi}{2}
\tightlist
\item
\end{enumerate}

\[
\begin{aligned}
g(x) = \ln f(x) =\left(1+\frac{1}{x} \right) \cdot \ln (1+x) = \frac{(1+x)\ln(1+x) }{x}\\
g'(x) = \frac{x+x\ln(1+x)-(1+x)\ln(1+x)}{x^{2}}=\frac{x-\ln(1+x)}{x^{2}}=\frac{f'(x)}{f(x)}\\
f'(x)=\left( \frac{x-\ln(1+x)}{x^{2}} \right)(1+x)^{\frac{1}{x}+1}
\end{aligned}
\]

\section{MATH-17}\label{math-17}

\textbf{定义1}

设 \(I \subset \mathbb{ R}\) 为区间，\(f:I\to \mathbb{R}\) 如果

\[
f((1-\lambda)x+\lambda y)\leq(1-\lambda) f(x)+\lambda f(y), \quad \forall x,y\in I, \lambda \in [0,1]
\]

称 \(f\) 为\textbf{凸函数}（\textbf{下凸函数}）

\(z=(1-\lambda)x+\lambda y, \; \lambda \in (0, 1)\) 分点公式

\(1-\lambda=\frac{y-z}{y-x}, \quad \lambda=\frac{z-x}{y-x}\)

\textbf{命题1}

设 \(f:I\to \mathbb{R}\) 为凸函数，\(x_{1},x_{2},x_{3} \in I\) 如果
\(x_{1}<x_{2}<x_{3}\)，则

\[
\frac{f(x_{2})-f(x_{1})}{x_{2}-x_{1}}\leq \frac{f(x_{3})-f(x_{1})}{x_{3}-x_{1}}\leq \frac{f(x_{3})-f(x_{2})}{x_{3}-x_{2}}
\]

\$\$

\begin{aligned}
x_{2}=\frac{x_{3}-x_{2}}{x_{3}-x_{1}}x_{1}+\frac{x_{2}-x_{1}}{x_{3}-x_{1}}x_{3} \\

f(x_{2})=f\left( \frac{x_{3}-x_{2}}{x_{3}-x_{1}}x_{1}+\frac{x_{2}-x_{1}}{x_{3}-x_{1}}x_{3} \right)
\leq \frac{x_{3}-x_{2}}{x_{3}-x_{1}}f(x_{1})+\frac{x_{2}-x_{1}}{x_{3}-x_{1}}f(x_{3}) \\

(x_{3}-x_{1})f(x_{2})\leq (x_{3}-x_{2})f(x_{1})+(x_{2}-x_{1})f(x_{3}) \\

\iff(x_{3}-x_{1})(f(x_{2})-f(x_{1})) \leq (x_{2}-x_{1})(f(x_{3})-f(x_{1})) \\

\iff (x_{3}-x_{1})(f(x_{2})-f(x_{3})) \leq (x_{3}-x_{2}) (f(x_{1})-f(x_{3}))


\end{aligned}

\$\$

\textbf{命题2}

设 \(f:I\to \mathbb{R}\) 则 \(f\) 为凸函数当且仅当
\(\forall x_{1},x_{2},x_{3}\in I, x_{1}<x_{2}<x_{3}\)

\[
\frac{f(x_{2})-f(x_{1})}{x_{2}-x_{1}}\leq \frac{f(x_{3})-f(x_{2})}{x_{3}-x_{2}}
\]

\$\$

\begin{aligned}
x_{1}=x, x_{3}=y, x_{2}=(1-\lambda)x_{1}+\lambda x_{3}\\
\lambda =\frac{x_{2}-x_{1}}{x_{3}-x_{1}} \in (0,1) \\

(x_{3}-x_{2})(f(x_{2})-f(x_{1})) \leq (x_{2}-x_{1}) (f(x_{3})-f(x_{2})) \\
\iff (x_{3}-x_{1})f(x_{2})\leq  (x_{3}-x_{2}) f(x_{1})+(x_{2}-x_{1})f(x_{3}) \\
\iff f((1-\lambda)x+\lambda y)=f(x_{2})\leq \frac{x_{3}-x_{2}}{x_{3}-x_{1}}f(x_{1})+\frac{x_{2}-x_{1}}{x_{3}-x_{1}}f(x_{3})= (1-\lambda)f(x)+\lambda f(y) \\
\forall x,y \in I, \; \lambda \in (0, 1), \; x\neq y
\end{aligned}

\$\$

\textbf{命题3}（弱极值定理）

设 \(a, b\in \mathbb{R},\; a<b, \; f:[a,b]\to \mathbb{R}\) 为凸函数，则

\[
\max f = \max \{  f(a), f(b) \}
\]

\[
\begin{aligned}
\forall t\in [a, b], \quad f(t)=f\left( \frac{b-t}{b-a}a+\frac{t-a}{b-a}b \right)\leq \frac{b-t}{b-a}f(a)+\frac{t-a}{b-a}f(b)\leq \max \{ f(a),f(b) \}
\end{aligned}
\]

\textbf{命题4}（强极值定理）

设 \(f:I\to \mathbb{R}\) 为凸函数，\(c \in I\)，\(c\) 不是 \(I\)
的端点，如果 \(f(c)=\max f\) 则 \(f \equiv f(c)\)

\[
\begin{aligned}
&x, y \in I, x<c<y \\ 
\implies &f(c)  \leq \frac{y-c}{y-x}f(x)  + \frac{c-x}{y-x}f(y)\leq f(c) \\
\implies &(y-x)f(c)=(y-c)f(x)+(c-x)f(y) \\ 
\implies &0\leq(y-c)(f(c)-f(x)) = (c-x)(f(y)-f(c))\leq 0 \\
\implies &f(x)=f(y)=f(c)
\end{aligned}
\]

\textbf{命题5}

设 \(f:(a,b)\to \mathbb{R}\) 为凸函数，\(x_{0}\in (a,b)\)，则
\(\exists k \in \mathbb{R},\text{ s.t. }\)

\[
f(x)\geq k(x-x_{0})+f(x_{0}), \quad \forall x \in (a, b)
\]

直线 \(y=k(x-x_{0})+f(x_{0})\) 叫凸函数的支撑线

证明

设
\(A= \left\{   \frac{f(x)-f(x_{0})}{x-x_{0}} \Bigm| x \in (a, x_{0}) \right\}\)

\(B=\left\{  \frac{f(x)-f(x_{0})}{x-x_{0}} \Bigm| x \in (x_{0}, b) \right\}\)

可得 \(x\leq y, \quad \forall x \in A, y\in B\)

则 \(\exists c, x\leq c\leq y, \quad \forall x \in A, y \in B\)

即
\(\frac{f(x)-f(x_{0})}{x-x_{0}}\leq c \implies f(x) \geq c(x-x_{0})+f(x_{0}), \quad \forall x \in (a,x_{0})\)

和
\(\frac{f(x)-f(x_{0})}{x-x_{0}} \geq c \implies f(x) \geq c(x-x_{0})+f(x_{0}), \quad \forall x \in (x_{0}, b)\)

又 \(f(x)=c(x-x_{0})+f(x_{0}), \quad x=x_{0}\)

综上得证

\section{MATH-18}\label{math-18}

\textbf{定义1}

设 \(f:A\subset \mathbb{R} \to \mathbb{R}, \; x_{0} \in A\)

如果 \(\forall \varepsilon  > 0 , \exists \delta >0 ,\text{ s.t. }\)

\[
|f(x)-f(x_{0})|<\varepsilon, \quad \forall x \in A, |x-x_{0}|<\delta
\]

称 \(f\) 在 \(x_{0}\) 点连续

如果 \(\forall x \in A\)，\(f\) 在 \(x\) 点连续，称 \(f\) 是连续函数或
\(f\) 连续

\textbf{例1}

\(f(x)=x, \;  x \in \mathbb{R}\) 则 \(f\) 连续

\textbf{例2}

\(f(x)=x^{2}, \;  x \in \mathbb{R}\) 证明 \(f\) 连续

证明

设 \(x_{0}\in \mathbb{R}\) 设 \(\varepsilon>0\) 令

\[
\delta = \frac{\varepsilon}{2|x_{0}|+1+\varepsilon}
\]

则 当 \(|x-x_{0}|<\delta\)

\[
|f(x)-f(x_{0})|\leq (|x|+|x_{0}|)|x-x_{0}|\leq(2|x_{0}|+1)\delta<\varepsilon
\]

\textbf{例3}

\$\$

H(x)=\left\{

\begin{aligned}
1, \quad x>0\\
0, \quad x\leq 0
\end{aligned}

\right. \$\$

不连续

证明

令 \(\varepsilon=\frac{1}{2}\) 设 \(\delta>0\)

令 \(x=\frac{\delta}{2}\) 则 \(x-0<\delta\)

\(|H(x)-H(0)|=1>\varepsilon\)

\textbf{命题1}

设 \(f:A\subset \mathbb{R}\to \mathbb{R}, \; x_{0} \in A\) 则 \(f\) 在
\(x_{0}\) 点连续当且仅当

\[
\begin{aligned}
&\forall \varepsilon>0 , \exists\delta>0,\text{ s.t.} \\
&|f(x)-f(x_{0})|\leq C\varepsilon, \forall x \in A, |x-x_{0}|<\delta 
\end{aligned}
\]

其中 \(C>0\) 为常数

\textbf{证明}

\(\implies\) 已知 \(f\) 满足定义，设 \(\varepsilon>0\) 则
\(C\varepsilon>0\) 则 \(\exists \delta>0,\text{ s.t.}\)

\[
|f(x)-f(x_{0})|<C\varepsilon, \quad \forall x \in A , \; |x-x_{0}|<\delta
\] \(\Longleftarrow\) 已知 \(f\) 满足命题1的条件，证 \(f\) 满足定义

设 \(\varepsilon>0\) 则 \(\frac{\varepsilon}{2C}>0\) 因此
\(\exists \delta>0,\text{ s.t.}\)

\[
|f(x)-f(x_{0})|\leq C \cdot \frac{\varepsilon}{2C} =\frac{\varepsilon}{2}<\varepsilon , \quad \forall x \in A, \; |x-x_{0}|<\delta
\]

\textbf{定义2}

设 \(f:A\subset \mathbb{R} \to \mathbb{R}, \; x_{0} \in A\)，如果
\(\exists \delta>0, \; M \in \mathbb{R}\) 使得

\[
|f(x)|\leq M, \quad \forall x \in A, |x-x_{0}|<\delta
\]

称 \(f\) 在 \(x_{0}\) 点局部有界

\textbf{命题2}

设 \(f:A\subset \mathbb{R}\to \mathbb{R}, \; x_{0}\in A\) 如果 \(f\) 在
\(x_{0}\) 点连续，则 \(f\) 在 \(x_{0}\) 点局部有界

\textbf{证明}

在定义1中，令 \(\varepsilon=1\)，则 \(\exists \delta >0,\text{ s.t.}\)

\[
|f(x)-f(x_{0})|\leq 1, \quad \forall x \in A,|x-x_{0}|<\delta
\]

则

\[
|f(x)|\leq |f(x_{0})|+1, \quad \forall x \in A, |x-x_{0}|<\delta 
\]

\textbf{命题3}

设 \(f,g:A\subset \mathbb{R} \to \mathbb{R}\) 在 \(x_{0}\in A\)
点连续，\(\alpha,\beta \in \mathbb{R}\)，则

\begin{enumerate}
\def\labelenumi{\arabic{enumi}.}
\tightlist
\item
  \(\alpha f+\beta g\) 在 \(x_{0}\) 点连续
\item
  \(fg\) 在 \(x_{0}\) 点连续
\item
  如果 \(\forall x \in A, g(x)\neq 0\) 则 \(\frac{f}{g}\) 在 \(x_{0}\)
  点连续
\end{enumerate}

\textbf{证明}

设 \(\varepsilon>0, \exists \delta_{1},\delta_{2}>0,\text{ s.t.}\)

\[
|f(x)-f(x_{0})|\leq\varepsilon, \quad \forall x \in A, |x-x_{0}|<\delta_{1}
\]

\[
|g(x)-g(x_{0})| \leq \varepsilon, \quad \forall x \in A, |x-x_{0}|< \delta_{2}
\]

\begin{enumerate}
\def\labelenumi{\arabic{enumi}.}
\tightlist
\item
\end{enumerate}

令 \(\delta=\min \{ \delta_{1},\delta_{2} \}\) 则当
\(x \in A, |x-x_{0}|<\delta\)

\[
|(\alpha f(x)+\beta g(x))-(\alpha f(x_{0})+\beta g(x_{0}))|\leq(|\alpha|+|\beta|+1)\varepsilon
\]

\begin{enumerate}
\def\labelenumi{\arabic{enumi}.}
\setcounter{enumi}{1}
\tightlist
\item
\end{enumerate}

\(g\) 在 \(x_{0}\)
点局部有界，\(\exists \delta_{3}>0, M \in \mathbb{R},\text{ s.t.}\)

\[
|g(x)|\leq M, \quad \forall x \in A, |x-x_{0}|<\delta_{3}
\]

令 \(\delta=\min\{ \delta_{1},\delta_{2},\delta_{3} \}\)

则当 \(x \in A, \; |x-x_{0}|<\delta\) 时

\[
\begin{aligned}
|f(x)g(x)-f(x_{0})g(x_{0})| &=|f(x)g(x)-f(x_{0})g(x)+f(x_{0})g(x)-f(x_{0})g(x_{0})| \\
&\leq |g(x)||f(x)-f(x_{0})|+|f(x_{0})||g(x)-g(x_{0})|\\
&\leq (M+|f(x_{0})|+1)\varepsilon
\end{aligned}
\]

\begin{enumerate}
\def\labelenumi{\arabic{enumi}.}
\setcounter{enumi}{2}
\tightlist
\item
\end{enumerate}

由定义1可得 \(\exists \delta_{3} > 0,\text{ s.t.}\)

\[
|g(x)-g(x_{0})|\leq \frac{|g(x_{0})|}{2} \implies |g(x)|\geq \frac{|g(x_{0})|}{2}
\]

令 \(\delta=\min\{ \delta_{1},\delta_{2},\delta_{3} \}\)

则 \(\forall x \in A, |x-x_{0}|>\delta\)

\[
\begin{aligned}
&\left|\frac{f(x)}{g(x)}-\frac{f(x_{0})}{g(x)}+\frac{f(x_{0})}{g(x)}-\frac{f(x_{0})}{g(x_{0})}\right| \\
= &\left| \frac{f(x)-f(x_{0})}{g(x)}+f(x_{0})\cdot \frac{g(x_{0})-g(x)}{g(x_{0})g(x)}\right|\\
\leq &  \left|  \frac{1}{g(x)}\right| \left(|f(x)-f(x_{0})|+\left| \frac{f(x_{0})}{g(x_{0})}\right| |g(x_{0})-g(x)|\right) \\
\leq & \left| \frac{2}{g(x_{0})}\right| \left(1+\left| \frac{f(x_{0})}{g(x_{0})}\right|\right) \varepsilon
\end{aligned}
\]

设 \(f,g:A\subset \mathbb{R}\to \mathbb{R}\) 在 \(x_{0} \in A\) 点连续

证明：

\begin{enumerate}
\def\labelenumi{\arabic{enumi}.}
\tightlist
\item
  \(|f|\) 在 \(x_{0}\) 点连续
\item
  \(\max\{ f,g \}, \min\{ f, g \}\) 在 \(x_{0}\) 点连续
\end{enumerate}

设 \(f:A_{1} \cup A_{2} \subset \mathbb{R}\to \mathbb{R}\)

其中 \(A_{1}\cap A_{2}=\varnothing\)

如果 \(f|_{A_{1}}, f|_{A_{2}}\) 连续，

\[
\inf \{ |x-y| \mid x \in A_{1}, y \in A_{2} \} > 0
\]

则 \(f\) 连续

\section{MATH-19}\label{math-19}

\textbf{定义}

设 \(f:A\subset \mathbb{R}\to \mathbb{R}, \; x_{0} \in A\) 如果
\(\forall \varepsilon>0, \exists \delta>0,\text{ s.t.}\)

\[
|f(x)-f(x_{0})|<\varepsilon, \quad \forall x \in A, |x-x_{0}|<\delta
\]

称 \(f\) 在 \(x_{0}\) 点连续

\textbf{命题}

设
\(f:A\subset \mathbb{R}\to \mathbb{R}, \; g:B\subset \mathbb{R}\to \mathbb{R}, \; \mathrm{Rg} f\subset B\)

设 \(x_{0}\in A, \; y_{0}=f(x_{0})\)，如果 \(f\) 在 \(x_{0}\)
点连续，\(g\) 在 \(y_{0}\) 点连续，则 \(g\circ f\) 在 \(x_{0}\) 点连续

\textbf{证明}

设 \(\varepsilon>0, \; \exists \sigma >0,\text{ s.t.}\)

\[
|g(y)-g(y_{0})|<\varepsilon, \quad \forall y \in B, |y-y_{0}|<\sigma
\]

\(\exists \delta>0,\text{ s.t.}\)

\[
|f(x)-f(x_{0})|<\sigma, \quad \forall x \in A, |x-x_{0}|<\delta
\]

设 \(x \in A, |x-x_{0}|<\delta\)

\[
|g(f(x))-g(f(x_{0}))|<\varepsilon
\]

初等函数都是连续的

\textbf{定理1}

定义在闭区间的函数一定是有界的

设 \(f:[a,b]\to \mathbb{R}\) 连续，则 \(f\) 有界

\textbf{证明} 反证法

假设 \(f\) 无界

记 \([a_{1},b_{1}]=[a,b]\)，令 \(c_{1}=\frac{a_{1}+b_{1}}{2}\)

如果 \(f\) 在 \([a_{1},c_{1}]\) 上无界，令
\([a_{2},b_{2}]=[a_{1},c_{1}]\)

否则，记 \([a_{2},b_{2}]=[c_{1},b_{1}]\)

重复上述操作

我们得到 \([a_{n},b_{n}],  \; n=1,2,\dots,\)

\begin{enumerate}
\def\labelenumi{\arabic{enumi}.}
\tightlist
\item
  \([a_{n},b_{n}]\supset [a_{n+1},b_{n+1}]\)
\item
  \(b_{n}-a_{n}=\frac{b-a}{2^{n-1}}\to 0\)
\item
  \(f\) 在 \([a_{n},b_{n}]\) 上无界
\end{enumerate}

令 \(A=\{ a_{n} \mid n=1,2,\dots, \}\)

\(B=\{ b_{n} \mid n=1,2,\dots, \}\)

则 \(x\leq y,\quad \forall x \in A, y \in B\)

\(\exists x_{0}\in \mathbb{R},\text{ s.t.}\)

\[
x\leq x_{0}\leq y, \quad \forall x \in A, y \in B
\]

在 \(x_{0}\) 点连续，所以在 \(x_{0}\) 点局部有界

\(\exists \varepsilon>0,\text{ s.t. } f\) 在
\([a, b]\cap(x_{0}-\varepsilon,x_{0}+\varepsilon)\) 上有界

\[
a_{n}\leq x_{0}\leq b_{n}, \quad \forall n=1,2,\dots,
\]

\(\exists N \in \mathbb{N}^{*},\text{ s.t.}\)

\[
b_{n}-a_{n}<\varepsilon
\]

所以
\([a_{n},b_{n}]\subset(x_{0}-\varepsilon,x_{0}+\varepsilon) \cap [a,b]\)

所以 \(f\) 在 \([a_{n},b_{n}]\) 上有界，与3矛盾

所以假设不成立，所以 \(f\) 有界

\textbf{定理2}

设 \(f:[a,b]\to \mathbb{R}\) 连续，则 \(f\) 有最大值（最小值）

\textbf{证明}

令 \(M=\sup f\) 下面证
\(\exists x_{0}\in[a,b],\text{ s.t. } f(x_{0})=M\)

记 \([a_{1},b_{1}]=[a,b]\) 令 \(c_{1}=\frac{a_{1}+b_{1}}{2}\)

\(\sup_{[a_{1},c_{1}]} f=M\) 或 \(\sup_{[c_{1},b_{1}]}f=M\)

若 \(\sup_{[a_{1},c_{1}]}f=M\) 则令 \([a_{2},b_{2}]=[a_{1},c_{1}]\)

否则 \([a_{2},b_{2}]=[c_{1},b_{1}]\)

重复上述操作

得到 \([a_{n},b_{n}],\;n=1,2,\dots,\)

\begin{enumerate}
\def\labelenumi{\arabic{enumi}.}
\tightlist
\item
  \([a_{n},b_{n}]\supset[a_{n+1},b_{n+1}]\)
\item
  \(b_{n}-a_{n}=\frac{b-a}{2^{n-1}}\to 0\)
\item
  \(\sup_{[a_{n},b_{n}]}f=M\)
\end{enumerate}

令 \(A=\{ a_{n} \mid n=1,2,\dots, \} \quad B=\{ b_{n}\mid n=1,2,.., \}\)

则 \(x\leq y, \; \forall x \in A, y\in B\)

\(\exists x_{0}\in \mathbb{R},\text{ s.t.}\)

\[
x\leq x_{0}\leq y, \quad \forall x \in A, y \in B
\]

即

\[
a_{n}\leq x_{0}\leq b_{n}, \quad n=1,2,\dots,
\]

下面证 \(f(x_{0})=M\) 反证，假设 \(f(x_{0})\neq M\) 则 \(f(x_{0})<M\)

\(\exists \varepsilon>0,\text{ s.t.}\)

\[
|f(x)-f(x_{0})|< \frac{M-f(x_{0})}{2}, \quad \forall x \in [a,b],|x-x_{0}|<\varepsilon
\]

则

\[
f(x)< \frac{f(x_{0})+M}{2}<M, \quad \forall x \in [a,b], |x-x_{0}|<\varepsilon
\]

所以在 \([a,b]\cap(x_{0}-\varepsilon,x_{0}+\varepsilon)\) 上
\(\sup f \leq \frac{f(x_{0})+M}{2}<M\)

\(\exists n \in \mathbb{N}^{*},\text{ s.t.}\)

\[
[a_{n},b_{n}]\subset[a,b]\cap(x_{0}-\varepsilon,x_{0}+\varepsilon)
\]

则 \(\sup_{[a_{n},b_{n}]}f\leq \frac{f(x_{0})+M}{2}<M\) 与3矛盾

所以假设不成立

所以 \(f(x_{0})=M\)

所以有最大值

\textbf{定理3}

设 \(f:[a,b]\to \mathbb{R}\) 连续，如果 \(f(a)f(b)\leq 0\) 则 \(f\)
必有零点

\section{MATH-21}\label{math-21}

\subsection{三角恒等式}\label{ux4e09ux89d2ux6052ux7b49ux5f0f}

\begin{enumerate}
\def\labelenumi{\arabic{enumi}.}
\tightlist
\item
  基本， \(\sin ^{2} x+\cos ^{2}x=1, \;\tan x= \frac{\sin x }{\cos x}\)
\item
  诱导公式， \(\sin(x+k\pi)=(-1)^{k}\sin x\)，
  \(\cos (x+k\pi)=(-1)^{k}\cos x\)，
  \(\sin\left( x+(2k+1)\frac{\pi}{2} \right) =(-1)^{k+1}\cos x\)，
  \(\cos\left( x+(2k+1) \frac{\pi}{2} \right)=(-1)^{k}\sin x\)
\item
  \(\sin(\alpha +\beta)=\sin \alpha \cos \beta+\cos \alpha \sin \beta\)，
  \(\sin(\alpha -\beta)=\sin \alpha \cos \beta-\cos \alpha \sin \beta\)，
  \(\cos(\alpha+\beta)=\cos \alpha \cos \beta-\sin \alpha \sin \beta\)，
  \(\cos(\alpha-\beta)=\cos \alpha \cos \beta+\sin \alpha \sin \beta\)，
  \(\tan(\alpha+\beta)=\frac{\tan \alpha+\tan \beta}{1-\tan \alpha \tan \beta}\)，
  \(\tan(\alpha-\beta )=\frac{\tan \alpha-\tan \beta}{1+\tan \alpha \tan \beta}\)
\item
  倍角公式， \(\sin(2\alpha)=2\sin \alpha \cos \alpha\)，
  \(\cos(2\alpha)=\cos ^{2}\alpha-\sin ^{2}\alpha=2\cos ^{2}\alpha-1=1-2\sin ^{2}\alpha\)，
  \(\tan(2\alpha)=\frac{2\tan \alpha}{1-\tan ^{2}\alpha}\)
\item
  降幂公式， \(\cos ^{2}\alpha=\frac{1+\cos 2\alpha}{2}\)，
  \(\sin ^{2}\alpha=\frac{1-\cos 2\alpha}{2}\)
\item
  半角公式，
  \(\sin \frac{\alpha}{2}=\pm \sqrt{ \frac{1-\cos \alpha}{2} }\)，
  \(\cos \frac{\alpha}{2}=\pm \sqrt{ \frac{1+\cos \alpha}{2} }\)，
  \(\tan \frac{\alpha}{2}=\pm \sqrt{ \frac{1-\cos \alpha}{1+\cos \alpha} }=\frac{\sin \alpha}{1+\cos \alpha}= \frac{1-\cos \alpha}{\sin \alpha}\)
\item
  万能公式，令 \(t=\tan \frac{\alpha}{2}\)，
  \(\sin \alpha=\frac{2t}{1+t^{2}}\)，
  \(\cos \alpha=\frac{1-t^{2}}{1+t^{2}}\)，
  \(\tan \alpha=\frac{2t}{1-t^{2}}\)
\item
  积化和差，
  \(2\sin \alpha \sin \beta=\cos(\alpha-\beta)-\cos(\alpha+\beta)\)，
  \(2\cos \alpha \cos \beta=\cos(\alpha-\beta)+\cos(\alpha+\beta)\)，
  \(2\sin \alpha \cos \beta=\sin(\alpha+\beta)+\sin(\alpha-\beta)\)，
  \(2\cos \alpha \sin \beta=\sin(\alpha+\beta)-\sin (\alpha-\beta)\)
\item
  和差化积，
  \(\cos x- \cos y = 2 \sin \left(  \frac{x+y}{2} \right)\sin\left( \frac{y-x}{2} \right)\)，
  \(\cos x+\cos y=2\cos\left( \frac{x+y}{2} \right)\cos\left( \frac{y-x}{2} \right)\)，
  \(\sin x+\sin y=2\sin\left( \frac{x+y}{2} \right)\cos\left( \frac{x-y}{2} \right)\)，
  \(\sin x-\sin y=2\cos\left( \frac{x+y}{2} \right)\sin\left( \frac{x-y}{2} \right)\)
\end{enumerate}

欧拉公式：\(e^{i\theta}=\cos \theta+\mathrm{i}\sin \theta\)

\textbf{练习}

证明：

\[
\frac{1}{2}+\sum_{k=1}^{n}\cos kx= \frac{\sin\left( n+\frac{1}{2} \right)x}{2\sin \frac{1}{2}x}
\]

\[
\sum_{k=0}^{n-1} \frac{ \sin\left( k+\frac{1}{2} \right)x }{\sin \frac{1}{2} x} = \frac{1-\cos nx}{1-\cos x}
\]

推导 \(\sin 3 \theta, \; \cos 3 \theta\) 的公式

\section{MATH-22}\label{math-22}

\begin{enumerate}
\def\labelenumi{\arabic{enumi}.}
\tightlist
\item
  \(\sin 3\theta=4\sin(60 ^ \circ -\theta)\sin \theta \sin (60^\circ+\theta)\)
\item
  \(\cos 3\theta=4\cos(60^{\circ}+\theta)\cos \theta \cos(60^{\circ}+\theta)\)
\item
  \(\sum_{i=1}^{n}\sin(\alpha+i\beta)\)
\item
  \(\sum_{i=1}^{n}\cos(\alpha+i\beta)\)
\end{enumerate}

\section{MATH-23}\label{math-23}

\section{MATH-24}\label{math-24}

\section{MATH-25}\label{math-25}

\section{MATH-26}\label{math-26}

\section{MATH-27}\label{math-27}

\subsection{复数}\label{ux590dux6570}

\(z=a+b \mathrm{i}, \; a, b \in \mathbb{R}\)

其中 \(\mathrm{i}=\sqrt{ -1 }\)，虚数单位

\(\mathrm{Re}z=a\)，实部；\(\mathrm{Im}z=b\)，虚部

\(b=0\)，实数；\(b\neq 0\)，虚数；\(a=0 \wedge b\neq 0\)，纯虚数

\(\mathbb{C}=\{ a+b\mathrm{i} \mid a, b \in \mathbb{R} \}\)，复数域

几何表示：复平面

\(z=a+b\mathrm{i} \iff\) 点 \((a,b) \iff\) 向量 \(\overrightarrow{OZ}\)

\begin{itemize}
\tightlist
\item
  共轭，\(\bar{z}=a-b\mathrm{i}\)
\item
  模，\(|z|=\sqrt{ a^{2}+b^{2} }\)
\item
  \(z_{1}+z_{2}=(a_{1}+a_{2})+\mathrm{i}(b_{1}+b_{2})\)
\item
  \(z_{1}z_{2}=(a_{1}a_{2}-b_{1}b_{2})+\mathrm{i}(a_{1}b_{2}+a_{2}b_{1})\)
\item
  \(\frac{z_{1}}{z_{2}}= \frac{z_{1}\overline{z_{2}}}{|z_{2}|^{2}}\)
\end{itemize}

共轭

\begin{itemize}
\tightlist
\item
  \(\bar{\bar{z}}=z\)
\item
  \(\overline{z_{1}\pm z_{2}}=\overline{z_{1}}\pm\overline{z_{2}}\)
\item
  \(\overline{z_{1}z_{2}}=\overline{z_{1}}\cdot\overline{z_{2}}\)
\item
  \(\overline{\left( \frac{z_{1}}{z_{2}} \right)}=\frac{\overline{z_{1}}}{\overline{z_{2}}}\)
\end{itemize}

模

\begin{itemize}
\tightlist
\item
  \(|z_{1}z_{2}|=|z_{1}||z_{2}|\)
\item
  \(| \frac{z_{1}}{z_{2}} |=\frac{|z_{1}|}{|z_{2}|}, \; z_{2}\neq 0\)
\item
  \(|z_{1}+z_{2}|\leq|z_{1}|+|z_{2}|\)
\end{itemize}

共轭与模

\begin{itemize}
\tightlist
\item
  \(|\overline{z}|=|z|\)
\item
  \(|z|=\sqrt{ z\overline{z} }\)
\end{itemize}

幅角

\(z \in \mathbb{C}\setminus \{  0 \}, \; r=|z|\)

\[
\frac{z}{r}=\cos \theta+\mathrm{i}\sin \theta
\]

\(\theta\)，\(z\) 的幅角；\(\theta \in [0, 2\pi)\)，\(z\)
的幅角主值，\(\theta=\mathrm{arg}z\)

\(\mathrm{Arg} z=\{ \mathrm{arg}z+2k\pi \mid k \in \mathbb{Z}\}\)

\textbf{命题1}

设 \(z_{1},z_{2},z_{3}\in \mathbb{C} \setminus \{ 0 \}\)

\begin{itemize}
\tightlist
\item
  \(-\mathrm{arg}z\in \mathrm{Arg}\overline{z}\)
\item
  \(-\mathrm{arg}z\in \mathrm{Arg} \frac{1}{z}\)
\item
  \(\mathrm{arg}z_{1}+\mathrm{arg}z_{2}\in \mathrm{Arg}(z_{1}z_{2})\)
\item
  \(\mathrm{arg}z_{1}-\mathrm{arg}z_{2}\in \mathrm{Arg}\left( \frac{z_{1}}{z_{2}} \right)\)
\end{itemize}

\(\theta_{1} = \mathrm{arg}z_{1}, \;\theta_{2}=\mathrm{arg}z_{2}\)

\(z_{1}=|z_{1}|(\cos \theta_{1}+\mathrm{i}\sin \theta_{1}),\;z_{2}=|z_{2}|(\cos \theta_{1}+\mathrm{i}\sin \theta_{2})\)

\(z_{1}z_{2}=|z_{1}||z_{2}|(\cos(\theta_{1}+\theta_{2})+\mathrm{i}\sin(\theta_{1}+\theta_{2}))\)

\textbf{命题2}

设 \(\theta_{1}, \theta_{2} \in \mathbb{R}\) 则

\(e^{\mathrm{i}\theta_{1}}\cdot e^{\mathrm{i}\theta_{2}}=e^{\mathrm{i}(\theta_{1}+\theta_{2})}\)

\(z^{n}\)，\(z^{0}=1\;(z\neq 0)\)，\(z^{-n}=\frac{1}{z^{n}}\; (z \neq 0)\)，\(n\in \mathbb{N}^{*}\)

\textbf{命题3}

设 \(n\in Z, \theta \in \mathbb{R}\) 则

\(\left(e^{\mathrm{i}\theta}\right)^{n}=e^{\mathrm{i}(n\theta)}\)

\section{MATH-28}\label{math-28}

设 \(z_{1},z_{2},z_{3},z_{4}\) 是复平面上的四个不同的点，称

\[
[z_{1},z_{2};z_{3},z_{4}]= \frac{z_{1}-z_{3}}{z_{2}-z_{3}} : \frac{z_{1}-z_{4}}{z_{2}-z_{4}}
\]

为 \(z_{1}, z_{2},z_{3},z_{4}\) 的\textbf{交比}

\textbf{定理}1

设 \(z_{1},z_{2},z_{3},z_{4}\) 是复平面上的四个不同的点，则

\begin{enumerate}
\def\labelenumi{\arabic{enumi}.}
\tightlist
\item
  \(z_{1},z_{2},z_{3},z_{4}\) 四点共线
  \(\iff [z_{1},z_{2};z_{3},z_{4}] \in \mathbb{R}, \frac{z_{1}-z_{3}}{z_{2}-z_{3}} \in \mathbb{R}\)
\item
  \(z_{1},z_{2},z_{3},z_{4}\) 四点共圆
  \(\iff [z_{1},z_{2};z_{3},z_{4}]\in \mathbb{R}, \frac{z_{1}-z_{3}}{z_{2}-z_{3}} \not\in \mathbb{R}\)
\end{enumerate}

\section{MATH-29}\label{math-29}

\textbf{定义}

\(f(z)= \frac{az+b}{cz+d}\) 称为分式线形变换

直线：\(\alpha x+\beta y+\gamma=0, \quad \alpha, \beta,\gamma \in \mathbb{R}, \alpha^{2}+\beta^{2}\neq 0\)

直线复数表示：\(\overline{B} z+B\bar{z}+C=0, \quad B \in \mathbb{C}, B\neq 0, C\in \mathbb{R}\)

圆的标准方程：\(|z-z_{0}|=r\)

\((z-z_{0})\overline{(z-z_{0})}=r^{2} \iff z\bar{z}-\overline{z_{0}}z-z_{0}\bar{z}+|z_{0}|^{2}-r^{2}=0\)

圆的一般方程：\(|z|^{2}+\overline{B}z+B\bar{z}+C=0, \quad B \in \mathbb{C},C\in \mathbb{R},|B|^{2}-C>0\)

统一形式：\(A|z|^{2}+\overline{B}z+B\bar{z}+C=0, \quad A,C\in \mathbb{R},B\in \mathbb{C},|B|^{2}-AC>0\)

\(\mathbb{C} \cup \{ \infty \}\) ：扩充复数域

在复平面中引入无穷远点，对应 \(\infty\)，叫扩充复平面

\begin{itemize}
\tightlist
\item
  \(z+\infty=\infty+z=\infty\)
\item
  \(z \cdot \infty=\infty \cdot z=\infty, \; z\neq 0\)
\item
  \(\frac{z}{\infty}=0\)
\item
  \(\frac{z}{0}=\infty, \; z\neq 0\)
\end{itemize}

统称直线和圆为圆；直线为过无穷远点的圆

设 \(b \in \mathbb{C}\)

\(T_{b}(z)=z+b, \; T_{b}^{-1}=T_{-b}\)

\(D_{r}(z)=rz, \; D_{r}^{-1}=D_{\frac{1}{r}}\)

\(R_{\theta}(z)=e^{i\theta}z, \; R_{\theta}^{-1}=R_{-\theta}\)

线形变换 \(f(z)=az+b, \; f=T_{b} \circ D_{|a|} \circ R_{\mathrm{arg}a}\)

\textbf{命题1}

设 \(f\) 是线形变换，则

\begin{enumerate}
\def\labelenumi{\arabic{enumi}.}
\tightlist
\item
  \([f(z_{1}),f(z_{2});f(z_{3}),f(z_{4})]=[z_{1},z_{2};z_{3},z_{4}]\)
\item
  如果 \(C\) 是圆，则 \(f(C)\) 也是圆
\end{enumerate}

圆 \(C: |z-z_{0}|=R\)

\(z \in \mathbb{C}, z\neq{}z_{0}\)

\[
z^{*}=\frac{R^{2}}{\overline{z-z_{0}}}+z_{0}
\]

\(z-z_{0}=re^{i\theta}, \; z^{*}-z_{0}=\frac{R^{2}}{r}e^{i\theta}\)

\(z^{*}\) 在射线 \(z_{0}z\) 上

\(z^{*}\) 称为 \(z\) 关于圆周 \(C\) 的对称点

\(z^{**}=z\)

\[
f(z)=\frac{R^{2}}{\overline{z-z_{0}}}+z_{0},\quad f^{-1}=f
\]

对称变换 \(S\)

\(S(z)=\frac{1}{z}, \; S^{-1}=S\)

\(S(0)=\infty, \; S(\infty)=0\)

\(S(z)\) 为 \(\bar{z}\) 关于单位圆的对称点

\textbf{命题2}

设 \(f(z)=\frac{1}{z}\)，则

\begin{enumerate}
\def\labelenumi{\arabic{enumi}.}
\tightlist
\item
  \([f(z_{1}),f(z_{2});f(z_{3}),f(z_{4})]=[z_{1},z_{2};z_{3},z_{4}]\)
\item
  如果 \(C\) 是圆，则 \(f(C)\) 也是圆
\end{enumerate}

分式线形变换

\[
f(z)=\frac{az+b}{cz+d}= \frac{a}{c}+\frac{b-\frac{ad}{c} }{cz+d}=\frac{a}{c}+\frac{bc-ad}{c^{2}z+cd}, \quad a, b, c, d \in \mathbb{C},ad-bc\neq{}0
\]

\[
f^{-1}(z)= \frac{-dz+b}{cz-a}
\]

\textbf{命题3}

设 \(f\) 是分式线形变换，则

\begin{enumerate}
\def\labelenumi{\arabic{enumi}.}
\tightlist
\item
  \([f(z_{1}),f(z_{2});f(z_{3}),f(z_{4})]=[z_{1},z_{2};z_{3},z_{4}]\)
\item
  如果 \(C\) 是圆，则 \(f(C)\) 也是圆
\end{enumerate}

\section{MATH-30}\label{math-30}
